\medterm{LRINEC Score} A blood-test score that estimates the chance of a serious soft-tissue infection such as necrotizing fasciitis. It uses inflammation, blood count, sodium, kidney, and blood-sugar results. The score cannot by itself confirm or exclude the infection; severe pain, rapid swelling, or serious illness may occur.

\synonyms
Laboratory Risk Indicator for Necrotizing Fasciitis; LRINEC

\medterm{Nabothian Cyst} A small, harmless mucus-filled lump on the cervix caused by a blocked cervical gland. It is common and usually causes no symptoms or need for treatment; an unusual or bleeding lesion still needs examination.

\medterm{Nabumetone} A nonsteroidal anti-inflammatory medicine used for pain and inflammation in selected joint or muscle conditions. It can irritate the stomach, affect kidneys, raise cardiovascular risk, or cause allergic reactions, especially with prolonged use.

\medterm{N-Acetylcysteine} An antidote used mainly for acetaminophen overdose. It helps protect the liver and may still help when treatment starts late or liver injury has already begun.

\medterm{Nadir} The lowest level reached by a measurement or symptom. After chemotherapy, it often means the period when blood-cell counts are lowest and the risk of infection or bleeding is greatest.

\medterm{Naegele Rule} A method for estimating a baby’s due date by adding seven days and about nine months to the first day of the last menstrual period. It is less reliable with irregular cycles or uncertain dates.

\medterm{Naegleria Fowleri} A free-living microscopic amoeba found mainly in warm freshwater lakes, rivers, and hot springs. If water containing it enters the nose, the amoeba can travel to the brain and cause primary amebic meningoencephalitis (PAM), a rare, rapidly worsening brain infection that is nearly always fatal. Early symptoms may include headache, fever, nausea, or vomiting, followed by stiff neck, confusion, seizures, or coma. Infection does not occur from swallowing the water and does not spread from person to person.

\medterm{Naegleria Infection} A rare infection caused when warm freshwater containing \textit{Naegleria fowleri} enters the nose and the amoeba reaches the brain. It causes primary amebic meningoencephalitis (PAM), which rapidly inflames and destroys brain tissue. Symptoms may include severe headache, fever, nausea, vomiting, stiff neck, confusion, seizures, or coma. It is not acquired by swallowing contaminated water and does not spread between people.

\medterm{Naevus} The British spelling of nevus, a usually benign localized skin or mucosal lesion such as a mole.

\medterm{Naevus Flammeus} A congenital vascular capillary malformation (port-wine stain) presenting as a flat, sharply demarcated, pink-to-violaceous cutaneous patch that does not regress spontaneously, classically involving the ophthalmic (V1) distribution of the trigeminal nerve in Sturge-Weber syndrome.
\synonyms
Port-Wine Stain; Capillary Vascular Malformation

\medterm{Naffziger Sign} A clinical physical sign for spinal nerve root compression or spinal space-occupying lesions; manual bilateral compression of the internal jugular veins causes transient elevation of intraspinal venous and CSF pressure, reproducing or worsening radiating radicular pain in the affected dermatome.
\synonyms
Naffziger's Sign; Jugular Compression Radiculopathy Test

\medterm{NAFLD} Abbreviation; see \textit{Non-Alcoholic Fatty Liver Disease}.

\medterm{Nail} A hard protective plate on the end of each finger and toe. Nails protect the tips, help with fine tasks, and may show changes caused by injury, infection, skin disease, or illness elsewhere.

\medterm{Nail Abnormalities} Changes in nail color, shape, thickness, texture, or attachment caused by trauma, infection, inflammatory disease, systemic illness, medication, or congenital conditions.

\medterm{Nail Changes} Alterations in nail color, shape, texture, or attachment that may result from local injury, skin disease, infection,
medicines, nutritional deficiency, or systemic illness. Their significance depends on
the pattern and other clinical findings.

\medterm{Nail Clipping For Fungal Testing} Laboratory examination of clipped nail material by microscopy, culture, or another validated test for fungal infection. Testing can help distinguish nail fungus from injury, psoriasis, or other nail disorders, especially before oral antifungal treatment.

\medterm{Nail Clubbing} Enlargement and rounding of the ends of the fingers or toes with more curved nails. It may run in families or signal disease of the lungs, heart, digestive system, or another organ, especially when it develops recently.

\medterm{Nail-Fold Capillaroscopy} A test that uses magnification to examine the tiny blood vessels near the nails. It can help distinguish primary Raynaud phenomenon from Raynaud phenomenon linked to conditions such as systemic sclerosis, dermatomyositis, or mixed connective-tissue disease.
\synonyms
Periungual Capillaroscopy; NFC Examination

\medterm{Nail Fungus} A fungal infection of a fingernail or toenail causing thickening, discoloration, brittleness, distortion, or separation from the nail bed.

\synonyms
Toenail Fungus

\medterm{Nail Furrows} Nail furrows are grooves or depressions running across or along a nail; transverse Beau lines can follow systemic illness or trauma, while longitudinal ridges are often age-related but may reflect local or systemic disease.

\medterm{Nail Infection} Infection of the nail or surrounding tissue, commonly caused by fungi or bacteria. It may cause thickening,
discoloration, separation of the nail, pain, swelling, or drainage; diabetes and poor circulation
increase the risk of complications.

\medterm{Nail Injuries} Trauma to the nail plate, nail bed, nail folds, or fingertip that may cause bruising, laceration, avulsion, fracture, or later nail-growth abnormalities.

\medterm{Nail Patella Syndrome} Usually an autosomal-dominant disorder caused by a pathogenic variant in \textit{LMX1B}, characterized by abnormal or absent nails, small or absent kneecaps, and elbow or pelvic abnormalities such as iliac horns. Glaucoma and kidney disease with proteinuria may also occur.

\medterm{Nail-Patella Syndrome} A rare inherited condition affecting the nails, kneecaps, elbows, and pelvis. Some people also develop kidney disease or glaucoma. The kneecaps may be small, absent, or displaced, and lifelong monitoring may be needed.

\medterm{Nail Pits} Small dents in the surface of a nail. They are common in psoriasis but may also occur with hair-loss disorders, eczema, lichen planus, or other conditions affecting nail growth.

\medterm{Nail Pitting} Punctate, cupuliform depressions in the dorsal nail plate resulting from parakeratotic defective keratinization within the proximal nail matrix, characteristically seen in psoriasis, alopecia areata, and reactive arthritis.

\medterm{Nail Psoriasis} Psoriasis affecting the nails can cause small pits, separation from the nail bed, yellow or red-brown patches, thickening, crumbling, or ridges. It may occur with skin psoriasis or psoriatic arthritis and can be severe or newly developed.

\medterm{Nails} Hard plates of keratin covering the ends of fingers and toes. Changes in color, shape, thickness, or attachment can result from injury, infection, skin disease, nutritional problems, or other illness.

\medterm{Nalbuphine} An opioid pain medicine with mixed effects, used for moderate to severe pain under medical supervision.

\medterm{Naloxone} A medicine that blocks opioids and can quickly restore breathing during an opioid overdose. It comes as a nasal spray or injection. Emergency help is still needed because its effect may wear off before the opioid does.

\medterm{Naltrexone} A longer-acting medicine that blocks opioid effects and reduces alcohol-related cravings in some people. Opioid use during treatment can trigger withdrawal.

\medterm{Nanocolloid} Very small particles carrying a radioactive tracer. After injection, they can show how lymph drains and help locate the first lymph node receiving drainage from a tumor.

\medterm{Nanomedicine} The use of very small materials or devices in medicine to help diagnose, prevent, monitor, or treat disease. Examples include nanoparticle drug carriers, imaging agents, biosensors, and some vaccine or cancer treatments; their effects and safety depend on the material, size, dose, and how the body handles them.

\medterm{Napelline} A poisonous plant chemical related to aconite alkaloids. Exposure can disturb nerve signals and heart rhythm, causing tingling, weakness, vomiting, dangerous blood-pressure changes, or life-threatening cardiac poisoning.

\medterm{Naphazoline} A topical alpha-adrenergic agonist used as a nasal decongestant or ocular vasoconstrictor; prolonged use can worsen congestion.

\medterm{Naphthalene Poisoning} Harm from mothballs or other products containing naphthalene. It can destroy red blood cells, causing weakness, jaundice, dark urine, kidney injury, seizures, or coma.

\medterm{Naproxen} A nonsteroidal anti-inflammatory drug used for pain, fever, and inflammation.

\medterm{Naproxen Sodium Overdose} Taking too much naproxen can cause stomach pain, vomiting, bleeding, drowsiness, kidney injury, seizures, or breathing problems. Large or intentional overdoses can be life-threatening.

\medterm{Narcissistic Personality Disorder} A long-lasting pattern of an unusually strong need for admiration, a sense of special importance, and difficulty recognizing other people's feelings. It causes problems in relationships or daily life and is diagnosed by a trained mental-health professional after a full assessment.

\synonyms
NPD; Megalomania; Pathological Narcissism

\medterm{Narcolepsy} A chronic disorder that affects the brain’s control of sleep and wakefulness. It causes overwhelming daytime sleepiness and sudden, irresistible sleep episodes, even after enough nighttime sleep; nighttime sleep may also be broken. The boundary between sleep and wakefulness is less clear, so some features of rapid-eye-movement (REM) sleep may occur while awake or when falling asleep or waking. Symptoms may include cataplexy, a brief loss of muscle control triggered by strong emotions; sleep paralysis; and vivid dream-like experiences. Type 1 narcolepsy includes cataplexy and is often associated with low hypocretin (orexin), a brain chemical that helps maintain wakefulness; type 2 usually does not include cataplexy.

\medterm{NARP Syndrome} A rare inherited mitochondrial disorder that can cause nerve damage, muscle weakness, loss of coordination, seizures, and vision problems. Symptoms and severity vary widely.
\synonyms
NARP; Neuropathy, Ataxia, and Retinitis Pigmentosa

\medterm{Narrow Angle Glaucoma} A condition in which the drainage area inside the eye is unusually narrow, raising the risk of a sudden rise in eye pressure. Severe eye pain, blurred vision, halos, nausea, or vomiting may occur during an acute attack.

\medterm{Narrowband UVB Phototherapy} A treatment that uses carefully measured ultraviolet light for conditions such as psoriasis, vitiligo, and eczema. Treatment requires eye protection and monitoring of the skin.

\medterm{Narrow-Complex Tachycardia} A fast heart rhythm that usually begins in the upper chambers of the heart. It produces narrow-looking beats on an ECG and may cause palpitations, dizziness, chest discomfort, breathlessness, or fainting.

\medterm{Narrow Pulse Pressure} A pulse pressure (systolic minus diastolic blood pressure) of less than 30 mmHg or less than twenty-five percent of systolic pressure. It reflects reduced stroke volume or marked peripheral vasoconstriction, commonly seen in severe aortic stenosis, cardiac tamponade, hypovolaemic shock, and advanced heart failure.

\medterm{Narrow Therapeutic Index} A medicine for which the helpful dose is close to the harmful dose. It requires careful dosing and sometimes blood tests to avoid toxicity.

\medterm{Nasal} Relating to the nose or the space inside it. The nose warms, moistens, and filters incoming air and helps with smell and speech. Nasal disease may cause blockage, bleeding, discharge, reduced smell, or difficulty breathing.

\medterm{Nasal Bleeding} Alternate name; see \textit{Epistaxis}.

\medterm{Nasal Cannula} A soft tube with two small prongs that rest inside the nostrils to deliver extra oxygen. It is used when a person needs mild to moderate breathing support, and the oxygen flow is adjusted while breathing is monitored.

\medterm{Nasal Cavity} The paired space inside the nose, extending from the nostrils to the nasopharynx. It warms, humidifies, and filters inspired
air and contains the receptors responsible for smell.

\medterm{Nasal Congestion} A blocked or stuffy feeling in the nose caused by swelling of the nasal lining, extra mucus, or a structural problem. It occurs with colds, allergies, sinus disease, irritants, and other conditions.

\medterm{Nasal Congestion And Rhinorrhea} Nasal obstruction or stuffiness together with a runny nose, commonly caused by viral infection, allergic or nonallergic rhinitis, sinus disease, irritants, or structural blockage. Persistent unilateral symptoms, recurrent bleeding, fever, facial swelling, or breathing difficulty requires evaluation.

\medterm{Nasal Corticosteroid Sprays} Nasal corticosteroid sprays reduce mucosal inflammation and congestion in allergic rhinitis, nasal polyps, and some chronic sinus disorders; correct direction away from the septum limits irritation and nosebleeds.

\medterm{Nasal Disorders} Conditions affecting the nose or nasal passages, including infection, allergy, trauma, obstruction, polyps, and tumors.

\medterm{Nasal Endoscopy} An examination in which a thin camera is used to look inside the nose and upper throat. It can help find swelling, polyps, a foreign object, bleeding, or a growth.

\medterm{Nasal Flaring} Widening and outward movement of the nostrils during breathing. It is a visible sign of
increased work of breathing, particularly in infants and children, and may accompany
airway obstruction, bronchiolitis, pneumonia, asthma, or severe systemic illness.

\medterm{Nasal Implant} A medical device placed inside or near the nose to support weakened tissue, change nasal shape, or improve airflow. The implant may be temporary or permanent, and risks include infection, movement, scarring, or breathing problems.

\medterm{Nasal Mucosal Biopsy} Removal of a small piece of tissue from the nose for laboratory examination. It may help investigate persistent inflammation, infection, unusual deposits, or a growth; soreness, bleeding, and infection are possible complications.

\medterm{Nasal Obstruction} Blocked airflow through one or both nasal passages. Causes include congestion, allergy, infection, a
deviated septum, polyps, trauma, or a mass. Persistent one-sided obstruction or bleeding requires
medical evaluation.

\medterm{Nasal Passage} One of the air channels inside the nose. Its lining warms, moistens, and filters inhaled air before it reaches the throat; swelling, mucus, polyps, or a structural blockage can make breathing difficult.

\medterm{Nasal Polyps} Soft, usually noncancerous growths in the lining of the nose or sinuses. They can cause a blocked nose, runny nose, and reduced sense of smell.

\synonyms
Sinonasal Polyps; Nasal Polyposis

\medterm{Nasal Septal Hematoma} A collection of blood inside the wall dividing the nostrils, usually after a nose injury. It can block breathing and damage the nose if it is not drained promptly.

\medterm{Nasal Septum} The wall that divides the nasal cavity into right and left sides. It consists mainly of cartilage in front and bone behind
and helps support the nose and direct airflow.

\medterm{Nasal Septum Perforation} A hole through the cartilage or bone separating the nasal passages, which may cause crusting, nosebleeds, whistling, or obstruction and can result from trauma, surgery, intranasal drugs, or inflammatory disease.

\medterm{Nasal Spray} A liquid or mist delivered into the nose. Different sprays contain saline, corticosteroids, antihistamines, or decongestants; correct use and duration depend on the product.

\medterm{Nasal Stent} A small device placed inside the nose to keep an airway open or support healing after surgery. It may cause crusting, blockage, infection, bleeding, or pressure injury.

\medterm{Nasal Vestibulitis} Inflammation or infection of the nasal vestibule, commonly caused by Staphylococcus bacteria and producing tenderness, crusting, redness, or small pustules near the nostrils.

\medterm{Nasogastric Feeding Tube} A flexible tube passed through the nose into the stomach to deliver enteral nutrition or medicines, or to decompress or aspirate the stomach.

\medterm{Nasogastric Tube} A tube passed through the nose into the stomach to remove contents or deliver nutrition, fluids, or medicines when clinically indicated.

\medterm{Nasojejunal Feeding Tube} An enteral feeding catheter passed through the nose, pharynx, stomach, and pylorus into the proximal jejunum under endoscopic or fluoroscopic guidance, indicated for patients requiring enteral nutrition who cannot tolerate gastric feeding due to severe gastroparesis, recurrent aspiration, or acute pancreatitis.

\synonyms
NJ Tube; Post-Pyloric Feeding Tube

\medterm{Nasolacrimal Duct} A membranous canal enclosed within a bony canal that conveys tears from the lacrimal sac downward into the inferior meatus of the nasal cavity beneath the inferior nasal concha, guarded by the mucosal valve of Hasner.

\synonyms
Tear Duct

\medterm{Nasopharyngeal Airway} A soft tube inserted through the nose to maintain a passage between the nares and
pharynx. It may be useful in a spontaneously breathing patient with partial upper-airway
obstruction but is relatively contraindicated in suspected basal skull fracture or
significant nasal trauma.

\medterm{Nasopharyngeal Carcinoma} A cancer arising in the upper throat behind the nose. It may cause a blocked or bleeding nose, hearing problems, a neck lump, headache, or swallowing trouble; risk varies with genetics, infections, and geographic background.

\synonyms
NPC; Nasopharyngeal Cancer

\medterm{Nasopharyngeal Culture} A laboratory test in which a sample from behind the nose is placed in a growth medium to look for selected bacteria or fungi. A negative result does not exclude every infection, and other tests may be more sensitive.

\medterm{Nasopharynx} The upper part of the throat behind the nose and above the soft palate. It carries air, contains the openings of the tubes that connect to the middle ears, and includes nearby lymph tissue that helps defend against infection.

\medterm{Natal Cleft} The groove between the buttocks, also called the intergluteal or sacral cleft.

\medterm{Natal Teeth} Teeth present at birth, most often involving the lower central incisors. They may be
loose, interfere with feeding, or injure the infant's tongue and require dental
assessment.

\medterm{Natriuresis} Removal of sodium from the body in urine. Water usually leaves with the sodium, which can reduce fluid volume and blood pressure; the process may be increased by natural hormones or water-removing medicines.

\medterm{Natriuretic} Causing the kidneys to remove more sodium in urine. Water usually leaves with the sodium, which can lower blood volume and pressure. The term describes certain body hormones and some medicines; related blood tests can help assess heart failure.

\medterm{Natural Immunity} Protection that develops after a person’s immune system responds to an infection, or temporary protection passed to a newborn through the placenta or breast milk. It may be incomplete or fade and is not always safer than vaccination.

\medterm{Natural Killer Cell} A type of white blood cell that can recognize and destroy some virus-infected or abnormal cells without previous exposure to that specific target. It helps provide early immune defense and releases signals that guide other immune cells.

\synonyms
NK Cells; Large Granular Lymphocytes

\medterm{Nausea} An unpleasant feeling that you may vomit. It can occur with stomach or intestinal illness, pregnancy, medicines, migraine, inner-ear problems, metabolic disorders, or serious abdominal or neurologic disease. Persistent vomiting, severe pain, dehydration, blood, or confusion may signal a serious condition.

\medterm{Nausea And Vomiting} Nausea is the urge to vomit; vomiting is forceful emptying of stomach contents. Causes include infection, medicines, pregnancy, migraine, obstruction, and other illnesses.

\medterm{Nausea And Vomiting – Adults} Nausea and vomiting in adults can result from infection, stomach or bowel disease, medicines, pregnancy, migraine, metabolic illness, neurologic disease, or blockage. Severe pain, blood, confusion, dehydration, chest symptoms, or persistent symptoms may indicate a serious condition.

\medterm{Nausea And Vomiting During Early Pregnancy} Nausea and vomiting are common in early pregnancy, but persistent symptoms with inability to keep fluids down, weight loss, dehydration, abdominal pain, fever, blood, or altered mental status may indicate hyperemesis gravidarum or another condition and require medical evaluation.

\medterm{Navel} The navel, or umbilicus, is the scar left after the umbilical cord separates and marks the former connection between fetus and placenta.

\medterm{Navicular} A boat-shaped bone in the middle of the foot that helps support the arch. Injury to it can cause foot pain and may heal slowly because part of the bone has a limited blood supply.

\medterm{NCS} Usually an abbreviation for nerve-conduction study, a test that measures how quickly and strongly electrical signals travel through nerves. It can help assess nerve damage, pressure on a nerve, or some causes of numbness and weakness.

\medterm{NDM-1} An enzyme made by some bacteria that makes them resistant to many important antibiotics. Infections caused by these bacteria can be difficult to treat and require laboratory testing and careful infection control.

\medterm{Near-Drowning} Survival for at least 24 hours following suffocation or submersion in a liquid medium, characterized pathologically by alveolar surfactant washout, pulmonary edema, ventilation-perfusion mismatch, acute respiratory distress syndrome, and hypoxic-ischemic encephalopathy.

\synonyms
Non-Fatal Drowning; Submersion Injury

\medterm{Nearsightedness} A common focusing problem in which nearby objects are clear but distant objects look blurry because light focuses in front of the retina. Glasses, contact lenses, or selected procedures can improve vision.

\synonyms
Myopia; Short-Sightedness

\medterm{Nebulizer} A device that converts liquid medicine into an aerosol for inhalation through a mouthpiece or mask. Common types include jet nebulizers, ultrasonic nebulizers, and vibrating-mesh nebulizers; they differ in how the aerosol is generated and delivered.

\synonyms
Nebuliser; Atomizer

\medterm{NEC} Abbreviation; see \textit{Neuroendocrine Carcinoma}.

\medterm{Necator} A group of hookworms that live in the small intestine. They enter through the skin, mature in the gut, and feed on blood, potentially causing abdominal symptoms, iron-deficiency anemia, weakness, and poor growth.

\medterm{Necator Americanus} A hookworm species that infects humans through skin penetration by larvae and causes intestinal blood loss.

\medterm{Neck} The neck connects the head to the trunk and contains the cervical spine, airway, esophagus, thyroid, vessels, nerves, muscles, and lymph nodes.

\medterm{Neck And Back Pain} Pain in the neck, upper back, lower back, or surrounding tissues, commonly caused by strain, age-related changes, disc problems, or trauma. Weakness, numbness, walking difficulty, bowel or bladder changes, fever, cancer history, or major trauma may indicate a serious cause.

\medterm{Neck Cancer} Cancer arising in tissues of the neck, such as the throat, thyroid, salivary glands, lymph nodes, or other structures. Symptoms may include a persistent lump, pain, swallowing trouble, voice change, or unexplained weight loss.

\medterm{Neck Dissection} Surgical removal of cervical lymph nodes and sometimes surrounding tissues to treat or stage head and neck cancer.

\medterm{Neck Lump} A localized swelling or mass in the neck caused by enlarged lymph nodes, thyroid or salivary disease, cyst, infection, or tumor; persistent or concerning masses require evaluation.

\medterm{Neck Mass} A lump or swelling in the neck caused by an enlarged lymph node, congenital cyst, thyroid or salivary-gland lesion, infection, or tumor. A persistent, enlarging, firm, fixed, or unexplained mass—especially in an adult—can have an important underlying cause.

\medterm{Neck Pain} Pain or stiffness in the neck, commonly caused by muscle strain, poor posture, injury, arthritis, degenerative disc disease, or nerve compression. It may be localized or spread to the shoulders, arms, or head.

\synonyms
Cervicalgia; Cervical Pain; Neck Stiffness

\medterm{Neck Pain And Dizziness} Neck pain with dizziness may reflect muscle strain, migraine, inner-ear disease, medication, blood-pressure change, or less commonly vascular or neurologic disease. Sudden severe symptoms, weakness, speech or vision change, fainting, or inability to walk may occur with a serious condition.

\medterm{Neck Sprain} Alternate name; see \textit{Whiplash}.

\medterm{Neck Strain} Alternate name; see \textit{Whiplash}.

\medterm{Neck X-Ray} Plain radiographic imaging of the cervical spine or neck soft tissues used to evaluate alignment, fractures, degenerative change, foreign bodies, or selected airway abnormalities.

\medterm{Necrobiosis Lipoidica} A chronic inflammatory skin disorder characterized by yellow-brown, atrophic plaques,
usually on the shins. It is associated with diabetes but can occur without diabetes.

\medterm{Necrobiosis Lipoidica Diabeticorum} A chronic inflammatory skin disorder producing yellow-brown atrophic plaques with a raised border, most often on the shins and associated with diabetes, although it can occur without diabetes.

\medterm{Necrolytic Acral Erythema} A rare rash with dark, thickened, or blistered patches, usually on the feet and ankles. It is often linked to long-term hepatitis C but can also occur with low zinc or other metabolic problems. Testing is needed to find the cause.
\synonyms{NAE; Hepatitis C-Associated Necrolytic Erythema}

\medterm{Necrolytic Migratory Erythema} A recurrent, erosive, erythematous eruption that often affects periorificial and
intertriginous skin. It is classically associated with glucagonoma but may have other causes,
so evaluation should consider nutritional, hepatic, pancreatic, and metabolic disease.

\medterm{Necrosis} Necrosis is death of cells or tissue within a living body caused by ischemia, infection, toxins, trauma, inflammation, or other injury.

\medterm{Necrotizing Anterior Scleritis} A severe inflammation of the white outer coat of the eye that causes tissue damage and can threaten vision. It usually causes intense eye pain, redness, and light sensitivity.

\medterm{Necrotizing Enterocolitis} A serious illness in which inflammation damages the intestine, mainly in premature or medically fragile newborns. A swollen abdomen, feeding problems, green vomit, bloody stools, or unusual sleepiness may occur.

\medterm{Necrotizing Fasciitis} A life-threatening infection that rapidly destroys tissue beneath the skin. Early skin changes may be mild despite severe pain and deeper damage. Treatment usually includes surgery and antibiotics.

\medterm{Necrotizing Granuloma} A granulomatous inflammatory lesion containing central tissue necrosis, commonly associated with infections such as tuberculosis or fungal disease but requiring clinical and histologic correlation.

\medterm{Necrotizing Infections} Severe infections that cause tissue death through germs, toxins, blocked blood flow, or intense inflammation. Examples include necrotizing fasciitis, necrotizing pneumonia, and necrotizing enterocolitis. Rapidly worsening pain, swelling, skin changes, severe illness, or organ problems may occur.

\medterm{Necrotizing Pneumonia} Severe pneumonia in which infection causes necrosis and cavitation of lung tissue, often accompanied by sepsis, respiratory failure, or abscess formation.

\medterm{Necrotizing Sialometaplasia} A benign, self-limiting, ischemic inflammatory lesion of the salivary glands (most commonly involving the minor salivary glands of the hard palate) that clinically and histologically mimics mucoepidermoid carcinoma or squamous cell carcinoma.
\synonyms
Ischemic Sialadenitis; Benign Palatal Sialometaplasia

\medterm{Necrotizing Soft Tissue Infection} A rapidly progressive infection of the tissue beneath the skin or in muscle that causes severe pain, swelling, severe illness, and tissue death. Treatment may include surgery and antibiotics.

\medterm{Necrotizing Vasculitis} Inflammation of blood-vessel walls causing fibrinoid necrosis and potential vessel narrowing, thrombosis, hemorrhage, ischemia, or organ injury.

\medterm{Needle} A hollow, sharp tube used to inject medicine or remove blood, fluid, or tissue. Needles come in different lengths and widths; the appropriate size depends on the procedure and body site.

\medterm{Needle Biopsy} Removal of cells, fluid, or a tissue sample through a needle for cytologic or histopathologic examination. Fine-needle aspiration primarily collects cells, whereas core biopsy collects tissue architecture.

\medterm{Needle Decompression} An emergency procedure used to release trapped air from around a collapsed lung when a tension pneumothorax is causing severe breathing or circulation problems. A trained clinician places a needle or catheter through the chest, then usually inserts a chest tube for ongoing treatment. It is a temporary lifesaving measure, not a complete repair.

\synonyms
Needle Thoracostomy; Tension Pneumothorax Decompression; Emergency Chest Decompression

\medterm{Needle Driver} A surgical instrument with gripping jaws and handles used to hold and guide a suture needle through tissue. Its design helps control the needle while protecting the operator's fingers.

\synonyms
Needle Holder
Suture Needle Holder

\medterm{Needle Holder} Alternate name; see \textit{Surgical Needle Holder}.

\medterm{Needlestick Injury} A puncture from a needle or sharp object that may expose someone to blood or body fluids. It can create a risk of infection, including HIV or hepatitis, and may lead to testing or preventive treatment.

\medterm{Neer Sign} A physical examination finding used to evaluate for subacromial impingement syndrome. The examiner stabilizes the patient's scapula while passively elevating the internally rotated arm in full forward flexion; reproduction of anterior shoulder pain indicates impingement of the rotator cuff tendons against the coracoacromial arch.

\synonyms
Neer Impingement Test

\medterm{Neer Test} A clinical diagnostic physical examination maneuver for subacromial shoulder impingement; the examiner stabilizes the patient's scapula to prevent upward rotation and passively performs forced forward elevation and internal rotation of the fully extended arm, driving the greater tuberosity against the anteroinferior acromion to reproduce anterior shoulder pain.
\synonyms
Neer Impingement Sign; Neer Shoulder Test

\medterm{Negative Feedback} A control process in which a change triggers a response that reduces or reverses the original change. It helps keep body functions, such as temperature, blood sugar, and hormone levels, within a useful range.

\medterm{Negative Predictive Value} The chance that a person with a negative test result truly does not have the condition. It depends on test accuracy and how common or likely the condition was before testing.

\medterm{Negative Pressure Pulmonary Edema} A form of non-cardiogenic pulmonary edema provoked by vigorous, forceful inspiratory efforts against an acutely obstructed upper airway (such as post-extubation laryngospasm), generating extreme negative intrathoracic pressures that increase venous return and cause alveolar capillary stress failure.

\synonyms
NPPE; Post-Obstructive Pulmonary Edema

\medterm{Negative Pressure Ventilation} Breathing support that uses pressure around the chest to help the lungs expand, instead of pushing air directly through the airway. It is used rarely, including for some neuromuscular conditions.

\medterm{Negative-Strand RNA Virus} A virus whose genetic material cannot be read directly to make proteins. After entering a cell, it must first make a readable copy of its RNA before producing new virus particles.

\medterm{Negative Symptoms} A reduction or loss of usual emotional expression, speech, motivation, pleasure, or interest in other people. They can occur in schizophrenia and some other conditions and may interfere with work, relationships, and daily activities.

\medterm{Neglect} Failure to provide a child with necessary food, shelter, supervision, education, emotional care, or medical care, causing harm or risk of harm. The term can also describe ignoring one side of space after brain injury; the intended meaning depends on context.

\medterm{Neglected Tropical Disease} An infectious or parasitic disease that mainly affects populations in tropical or subtropical regions with limited resources. Examples include schistosomiasis, leprosy, dengue, and lymphatic filariasis.

\medterm{Neglect Syndrome} A neurologic problem in which a person fails to notice or respond to one side of the body or surrounding space, usually after injury to the opposite side of the brain. It is more than ordinary inattention.

\medterm{Neisseria Gonorrhoeae} A gram-negative diplococcus that causes gonorrhea, a sexually transmitted infection. The
organism has developed resistance to many antibiotics including fluoroquinolones and
macrolides.

\medterm{Neisseria Infection} Infection caused by a pathogenic Neisseria species. The most important human infections are gonorrhoea caused by N. gonorrhoeae and meningococcal disease caused by N. meningitidis; illness may affect the genital tract, throat, rectum, joints, bloodstream, or meninges, and treatment depends on the site and antibiotic susceptibility.

\synonyms
Neisseria Infections

\medterm{Neisseria Meningitidis} A gram-negative diplococcus that causes meningitis and meningococcemia. The organism has
a polysaccharide capsule that is the target of conjugate vaccines. Serogroups A, B, C,
W, X, and Y cause most disease. Meningitis presents with fever, headache, neck
stiffness, and altered mental status.

\medterm{Nelson Syndrome} Enlargement and hypersecretion of an adrenocorticotropic-hormone-producing pituitary
tumour after bilateral adrenalectomy, usually for Cushing disease. It may cause rapid
skin hyperpigmentation, headache, visual-field loss, and other effects of an expanding
sellar mass.

\medterm{Nemaline Myopathy} A group of inherited muscle diseases that cause weakness because muscle fibers do not develop or work normally. Severe forms begin in infancy and may affect movement, feeding, or breathing; milder forms begin later. The severity varies widely, and care may include breathing support, nutrition, physical therapy, and genetic counseling.

\synonyms
Rod Body Myopathy; Nemaline Rod Myopathy

\medterm{Neoadjuvant Therapy} Treatment given before the main treatment, such as surgery or radiotherapy, to shrink a tumor or make it easier to treat. It may also show how the tumor responds and help doctors plan the next treatment.

\medterm{Neocortex} The outer part of the cerebral hemispheres involved in conscious perception, language, planning, memory, and voluntary movement.

\medterm{Neoehrlichiosis} An infection caused by the bacterium Candidatus Neoehrlichia mikurensis, usually transmitted by ticks. It can cause fever and blood-vessel inflammation, especially in people with weakened immunity.

\medterm{Neologisms} Newly created words, phrases, or idiosyncratic combinations of word fragments that have meaning only to the speaker and are incomprehensible to others. They are characteristic features of formal thought disorder in schizophrenia and fluent (Wernicke) aphasia.

\medterm{Neonatal} Relating to a newborn, especially during the first 28 days after birth. This period includes important changes in breathing, feeding, temperature control, immunity, circulation, and adjustment to life outside the womb.

\medterm{Neonatal Abstinence Syndrome} Withdrawal symptoms in a newborn after exposure to opioids or certain other medicines or substances during pregnancy. Symptoms may include trembling, irritability, poor feeding, vomiting, diarrhea, sweating, or breathing problems.

\medterm{Neonatal Cephalohematoma} A collection of blood beneath the outer covering of a newborn’s skull, usually caused by birth trauma. It is limited by the skull’s seams and generally resolves over time, but it can be associated with injury or jaundice.

\medterm{Neonatal Conjunctivitis} Inflammation of the eye’s outer lining during the first month of life, caused by chemical irritation, bacteria, or viruses acquired during delivery. Pus-like discharge, eyelid swelling, and corneal involvement may occur.

\medterm{Neonatal Cystic Fibrosis Screening Test} A newborn screening test, usually measuring immunoreactive trypsinogen with reflex DNA or sweat testing, used to identify infants at increased risk for cystic fibrosis.

\medterm{Neonatal Diabetes Mellitus} Diabetes first diagnosed during the first six months of life. It may be temporary or lifelong and is usually caused by a genetic change affecting insulin production; specialist care is needed.

\medterm{Neonatal Galactorrhea} A temporary milk-like breast secretion in a newborn caused by exposure to maternal hormones. It usually resolves without treatment, and pressure on the breasts can cause irritation or infection.

\synonyms
Neonatal Milk Secretion; Newborn Galactorrhea

\medterm{Neonatal Herpes Simplex} Herpes simplex infection in a baby during the first 28 days of life, usually acquired around delivery. It may affect the skin, eyes, mouth, brain, or several organs and can resemble sepsis. Treatment uses antiviral medicine.

\medterm{Neonatal Herpes Simplex Virus} A severe neonatal infection caused by herpes simplex virus, typically acquired during
passage through the birth canal, presenting within the first month of life as skin, eye,
and mouth disease, central nervous system disease, or disseminated disease with
multiorgan involvement.

\medterm{Neonatal Hyperbilirubinemia} An excess of bilirubin in a newborn’s blood, causing yellow skin or eyes. Mild cases are common, but rapidly rising or very high levels can damage the brain; treatment may include frequent feeding, phototherapy, or rarely an exchange transfusion.

\medterm{Neonatal Hyperglycemia} An abnormally elevated blood glucose level in infants and neonates (typically defined as blood glucose $>145\text{ mg/dL}$ or $>8.0\text{ mmol/L}$). Common causes include extreme prematurity, intravenous dextrose overinfusion, sepsis, severe physiological stress, corticosteroid therapy, or transient neonatal diabetes mellitus.

\synonyms
Infantile Hyperglycemia; High Blood Glucose in Infants; Newborn Hyperglycemia

\medterm{Neonatal Hypertension} Sustained blood pressure elevation in infants and neonates exceeding age- and birthweight-adjusted norms. Frequently secondary to umbilical artery catheter-related renal artery thrombosis, congenital renal anomalies, bronchopulmonary dysplasia, or congenital heart defects.

\synonyms
Infant Hypertension; Neonatal High Blood Pressure; Infantile Hypertension

\medterm{Neonatal Hypocalcemia} Low serum calcium in a newborn, which may cause jitteriness, poor feeding, hypotonia, apnea, or seizures. Causes include prematurity, maternal diabetes, birth asphyxia, and disorders of parathyroid function.

\medterm{Neonatal Hypoglycemia} Blood sugar that is too low in a newborn. It may cause trembling, sleepiness, poor feeding, breathing pauses, or seizures and is more likely with prematurity, unusual birth size, or maternal diabetes; prompt monitoring and treatment are needed.

\medterm{Neonatal Hypothermia} A newborn’s body temperature below 36.5°C. It can result from exposure, prematurity, infection, or illness and may worsen breathing, blood sugar, and circulation problems. Prompt warming and assessment are needed; skin-to-skin contact and a warm, dry environment help prevent it.

\medterm{Neonatal Hypothyroidism} Insufficient thyroid hormone production in a newborn, which can impair brain development and growth if untreated; newborn screening enables early diagnosis and treatment.

\medterm{Neonatal Hypoxic-Ischemic Encephalopathy} Brain injury in a newborn caused by too little oxygen or blood flow around the time of birth. It may cause poor alertness, abnormal muscle tone, seizures, breathing difficulty, or injury to other organs. Selected babies benefit from carefully controlled cooling started soon after birth.

\medterm{Neonatal Infections} Infections occurring during the first weeks of a newborn’s life. They may begin before or soon after birth from exposure during pregnancy or delivery, or later from the environment or medical devices. Warning signs include poor feeding, temperature change, breathing difficulty, unusual sleepiness, irritability, or color change.

\medterm{Neonatal Jaundice} Yellow discoloration of a newborn's skin or eyes caused by elevated bilirubin. It is common but may
require treatment or investigation when severe, early, or prolonged.

\synonyms
Neonatal Hyperbilirubinemia; Icterus Neonatorum

\medterm{Neonatal Lupus} A condition caused by antibodies passed from a pregnant person to a baby. It may cause a temporary rash, low blood counts, liver problems, or a slow heart rhythm. Most effects improve as the antibodies leave the baby's blood, but heart block can be serious and may require a pacemaker.

\synonyms
Neonatal Lupus Erythematosus; NLE

\medterm{Neonatal Mechanical Ventilation} Invasive or non-invasive mechanical positive-pressure ventilation specifically tailored for neonates and infants with respiratory failure, severe respiratory distress syndrome (RDS), bronchopulmonary dysplasia, or apnea of prematurity. Emphasizes lung-protective strategies such as high-frequency oscillatory ventilation and synchronized intermittent mandatory ventilation.

\synonyms
Pediatric Mechanical Ventilation; Neonatal Ventilatory Support; Infant Mechanical Ventilation

\medterm{Neonatal Neutropenia} An absolute neutrophil count (ANC) below $1{,}000/\mu\text{L}$ to $1{,}500/\mu\text{L}$ in infants, or below $500/\mu\text{L}$ in severe cases. Etiologies include neonatal sepsis, maternal preeclampsia, alloimmune or autoimmune neutropenia, twin-to-twin transfusion, bone marrow failure syndromes (such as Kostmann syndrome), or drug exposure.

\synonyms
Infant Neutropenia; Neonatal Low Neutrophil Count; Pediatric Neutropenia

\medterm{Neonatal Opioid Withdrawal Syndrome} Withdrawal symptoms in a newborn exposed to opioid medicines or drugs during pregnancy. The baby may be hard to settle, tremble, feed poorly, vomit, have loose stools, sweat, or breathe rapidly. Quiet surroundings, swaddling, skin-to-skin contact, and feeding support are used first; some babies need hospital medicines and monitoring.

\synonyms
NOWS; Opioid-Induced NAS

\medterm{Neonatal Peripheral Arterial Line} A fine indwelling cannula placed in a peripheral artery (such as the radial, posterior tibial, or dorsalis pedis artery) of an infant for continuous invasive hemodynamic arterial blood pressure monitoring and frequent arterial blood gas sampling in neonatal intensive care.

\synonyms
Infant Arterial Line; Pediatric A-Line; Neonatal Arterial Cannulation

\medterm{Neonatal Peripheral IV Line} A short peripheral venous catheter inserted into a superficial vein in the hand, foot, forearm, or scalp of an infant for short-term hydration, medication delivery, or blood transfusions.

\synonyms
Infant PIV; Pediatric Peripheral Cannula; Neonatal IV Cannulation

\medterm{Neonatal PICC Line} A small-caliber (1.0--2.0 Fr) central venous catheter inserted percutaneously into a peripheral vein of an infant or neonate with the tip positioned in the superior or inferior vena cava, providing reliable vascular access for total parenteral nutrition and hyperosmolar medications.

\synonyms
Neonatal Percutaneous Central Line; Infant PICC; Pediatric Percutaneous Central Catheter

\medterm{Neonatal Pneumothorax} Pathologic accumulation of extrapulmonary air within the pleural space in a newborn or infant, which may occur spontaneously or secondary to meconium aspiration, respiratory distress syndrome, or mechanical barotrauma.

\synonyms
Infant Pneumothorax; Neonatal Air Leak Syndrome; Newborn Pleural Air Collection

\medterm{Neonatal Polycythemia} An abnormally high red-cell concentration in a newborn, often defined by a venous hematocrit of at least 65 percent. It may increase blood viscosity and cause respiratory, neurologic, or metabolic symptoms.

\medterm{Neonatal Respiratory Distress Syndrome} Breathing difficulty in a premature newborn caused mainly by insufficient pulmonary surfactant, leading to stiff lungs and impaired oxygen exchange.

\medterm{Neonatal Resuscitation} Emergency steps used to help a newborn who is not breathing well or has a slow heart rate after birth. They may include warmth, airway support, assisted breathing, and chest compressions.

\medterm{Neonatal Sepsis} A serious infection causing illness throughout a newborn's body during the first 28 days of life. It may begin before or after birth and can cause poor feeding, temperature change, breathing difficulty, sleepiness, or color change.

\medterm{Neonatal Skin Conditions} Skin changes common in newborns, including temporary rashes, tiny white spots, heat rash, and small yellow bumps. Most are harmless, but blistering, widespread redness, fever, poor feeding, or a rapidly worsening rash may signal illness.

\medterm{Neonatal Testicular Torsion} Twisting of the spermatic cord around a testicle during the newborn period, reducing or stopping its blood supply. It may appear as a firm, enlarged, discolored, or painless scrotal swelling.

\medterm{Neonatal Thyrotoxicosis} Excess thyroid activity in a newborn, usually caused by thyroid-stimulating antibodies passed from the mother. It may cause a fast heartbeat, irritability, poor feeding, poor weight gain, sweating, or breathing problems and usually needs temporary treatment and monitoring.

\medterm{Neonatal Transition} The newborn’s adjustment from receiving oxygen through the placenta to breathing with the lungs after birth. The lungs fill with air, blood flow changes, and the baby must begin controlling temperature, blood pressure, and feeding.

\medterm{Neonate} A neonate is a newborn infant during the first 28 completed days after birth, a period requiring attention to feeding, temperature, jaundice, infection, breathing, and transition to extrauterine life.

\medterm{Neonatologist} A pediatrician who specializes in the medical care of newborn infants, particularly
premature and critically ill neonates.

\medterm{Neoplasia} Abnormal, uncontrolled growth of cells that forms a new mass or tissue. It may be benign and remain localized, or malignant and invade nearby tissues or spread elsewhere.

\medterm{Neoplasm} A new, abnormal growth of cells that forms a mass or tissue. A neoplasm may be benign and stay localized, or malignant and invade nearby tissues or spread to distant organs.

\medterm{Neoplastic} Relating to an abnormal new growth of cells. A neoplastic growth may be noncancerous and localized or cancerous and able to invade nearby tissue or spread.

\medterm{Neostigmine} An acetylcholinesterase inhibitor used to reverse selected nondepolarizing neuromuscular blockade and to treat some disorders of neuromuscular transmission or bowel and bladder motility.

\medterm{Nephoptosis} Abnormal downward movement or mobility of a kidney, also called floating kidney.

\medterm{Nephrectomy} Surgery to remove part or all of a kidney. It may treat kidney cancer, severe kidney damage, infection, or a kidney that no longer works.

\medterm{Laparoscopic Nephrectomy} A minimally invasive surgical procedure in which a diseased or donor kidney is dissected and extracted using laparoscopy or robotic assistance, categorized as radical (en bloc removal with Gerota fascia and adrenal gland for renal cell carcinoma) or partial (nephron-sparing surgery).

\medterm{Nephric} Relating to the kidney. In embryology, nephric structures are involved in development of the urinary system.

\medterm{Nephritic Syndrome} A pattern of kidney inflammation that causes blood in the urine, reduced urine output, swelling, and high blood pressure. Protein may also leak into the urine. It can reduce kidney function and requires evaluation to find and treat the cause.

\medterm{Nephritis} Inflammation of kidney tissue that can cause blood or protein in the urine, swelling, high blood pressure, reduced urine, or loss of kidney function. Infection, medicines, and immune diseases are possible causes.

\medterm{Nephroblastoma} Another name for Wilms tumor, the most common kidney cancer in children. It may appear as a painless abdominal lump, abdominal pain, blood in the urine, fever, or high blood pressure and needs specialist treatment.

\medterm{Nephrocalcinosis} Buildup of calcium deposits inside kidney tissue. It may result from high calcium levels, hormone or kidney disorders, inherited conditions, dehydration, or medicines and can lead to stones, blood in the urine, or reduced kidney function.

\medterm{Nephrogenic Diabetes Insipidus} A condition in which the kidneys do not respond properly to vasopressin, the hormone that helps conserve water. It causes large amounts of dilute urine and intense thirst; inherited changes, kidney disease, electrolyte problems, or medicines may cause it.

\medterm{Nephrogenic Fibrosing Dermopathy} An older name for nephrogenic systemic fibrosis, a rare condition causing thick, tight, hardened skin and sometimes stiff joints in people with severe kidney disease, usually after exposure to certain gadolinium contrast agents.

\medterm{Nephrogenic Systemic Fibrosis} A rare disorder in which skin and connective tissue become thick and hardened, sometimes limiting movement. It is linked mainly to certain gadolinium contrast exposures in people with severe kidney impairment and may affect internal tissues.

\medterm{Nephrolithiasis} Formation of stones in the kidneys or urinary tract. Stones can cause severe side or back pain, blood in the urine, nausea, or blockage; dehydration, inherited conditions, medicines, and diet can contribute.

\medterm{Nephrologist} A doctor who specializes in kidney disease. Nephrologists assess reduced kidney function, high blood pressure related to kidney problems, kidney inflammation, and mineral or salt disturbances and may manage dialysis or transplant care.

\medterm{Nephrology} The medical specialty concerned with kidney disease, fluid and electrolyte disorders, hypertension, dialysis, and kidney transplantation.

\medterm{Nephrolysis} Surgical freeing of a kidney from surrounding adhesions or tissue.

\medterm{Nephron} The microscopic working unit of a kidney. Each nephron filters blood, returns needed water and substances to the body, and removes waste to make urine; each kidney contains about one million.

\medterm{Nephron Function} Each kidney contains about one million nephrons, its microscopic filtering units. A nephron filters blood, removes wastes into urine, and returns needed water and substances to the bloodstream while helping control salts, acid, and blood pressure.

\medterm{Nephronophthisis} A rare inherited kidney disease that gradually reduces the kidneys’ ability to concentrate urine. It may cause excessive thirst and urination, anemia, growth problems, and eventually kidney failure.

\medterm{Nephropathy} Disease or damage of the kidneys. It may reduce waste removal, disturb fluid and salt balance, cause protein or blood in the urine, or eventually lead to kidney failure.

\medterm{Nephroptosis} Abnormal downward movement of a kidney when a person stands, sometimes causing flank discomfort, urinary symptoms, or intermittent obstruction.

\synonyms
Wandering Kidney; Ren Mobilis; Hypermobile Kidney

\medterm{Nephrosclerosis} Hardening and scarring of kidney tissue, commonly caused by long-standing high blood pressure or blood-vessel disease. Progressive damage can reduce filtering ability, cause protein in urine, and lead to chronic kidney disease.

\medterm{Nephrostomy} A procedure that creates a drainage path from the kidney through the skin when urine cannot pass normally to the bladder. A temporary tube usually carries urine into a bag while the blockage is treated.

\medterm{Nephrostomy Tube} A drainage tube passed through the skin into the kidney to carry urine into a bag when the ureter is blocked. Blockage may be caused by a stone, tumor, or narrowing. Infection with blockage can require urgent drainage; the tube may be temporary or long-term.

\medterm{Nephrotic Range Proteinuria} Very large protein loss in the urine, usually more than 3.5 grams per day in an adult. It often indicates serious damage to the kidney filters and can cause swelling, low blood protein, high cholesterol, blood clots, or infection.

\medterm{Nephrotic Syndrome} A kidney disorder caused by damage to the kidney's tiny filtering units, allowing large amounts of vital protein (albumin) to spill into the urine. This loss leads to low blood protein, high cholesterol, frothy urine, and noticeable swelling (edema), especially around the eyes in the morning and in the feet and legs. It also increases the risks of infections and blood clots; common causes include minimal change disease in children and diabetes or kidney inflammation in adults.
\synonyms Nephrosis; Protein-losing nephropathy; Glomerular nephrosis.

\medterm{Nephrotoxic} Describing a substance or condition that can damage the kidneys or reduce their function. Examples include some medicines, poisons, infections, and severe dehydration; risk depends on dose, duration, and kidney health.

\medterm{Nephrotoxicity} Kidney damage caused by a medicine, poison, contrast material, or other substance. It may reduce urine output or raise blood creatinine; risk is lowered by avoiding triggers when possible and adjusting doses for kidney function.

\medterm{Nernst Equation} A fundamental biophysical formula that calculates the equilibrium electrical potential across a semipermeable cell membrane for a specific single ion species based on its intracellular and extracellular concentration gradient at physiological body temperature.

\medterm{Nerve} A bundle of nerve fibers that carries messages between the brain or spinal cord and the rest of the body. Nerves control movement and carry sensations such as touch, pain, and temperature.

\medterm{Nerve Biopsy} Removal of a small piece of a nerve for laboratory examination when serious nerve disease is suspected. It may help identify inflammation, infection, abnormal deposits, or a tumor, but can leave permanent numbness or weakness near the biopsy site.

\medterm{Nerve Block} An injection of numbing medicine near a nerve or group of nerves to temporarily stop pain signals. It may help diagnose the source of pain or provide pain relief during or after a procedure; weakness or other temporary effects can occur.

\medterm{Nerve Compression} Alternate name; see \textit{Pinched Nerve}.

\medterm{Nerve Conduction Studies} Tests that stimulate peripheral nerves with small electrical pulses and record how quickly and strongly signals travel. They help identify nerve damage, compression, loss of the nerve’s insulating covering, or loss of nerve fibers.

\medterm{Nerve Conduction Study} A test that gives small electrical pulses to a nerve and measures how quickly and strongly the signal travels. It helps identify nerve damage, pressure on a nerve, or loss of its protective covering.

\medterm{Nerve Conduction Velocity} The speed at which an electrical signal travels along a peripheral nerve. A nerve-conduction test measures this speed and the signal strength to help identify nerve damage, pressure on a nerve, or damage to its insulating covering.

\medterm{Nerve Deafness} Hearing loss caused by damage to the cochlea, auditory nerve, or central auditory pathways rather than blockage of sound transmission through the outer or middle ear.

\synonyms
Sensorineural Hearing Loss; SNHL

\medterm{Nerve Endings} The tiny ends of nerves that detect sensations such as touch, pain, heat, or pressure, or release signals to muscles and glands. Injury or irritation can cause pain, numbness, tingling, or abnormal sweating.

\medterm{Nerve Fiber Types} A physiological classification of peripheral nerve axons based on diameter and conduction velocity: group A (thickly myelinated motor and sensory fibers: alpha, beta, gamma, delta), group B (thinly myelinated preganglionic autonomic fibers), and group C (unmyelinated, slow postganglionic and pain/temperature C-fibers).

\synonyms
Erlanger-Gasser Classification

\medterm{Nerve Injury} Damage to a nerve that can cause pain, numbness, tingling, weakness, altered reflexes, or changes in sweating, blood pressure, digestion, or bladder function. Recovery depends on the nerve, severity, cause, and treatment.

\medterm{Nerve Sheath Myxoma} A benign myxoid peripheral-nerve-sheath tumor, usually presenting as a slow-growing papule or nodule on the fingers, hands, or upper extremities.

\medterm{Nerve Sparing} A surgical approach that tries to preserve nearby nerves needed for erection, sensation, or other functions. It may improve recovery after prostate or other pelvic surgery, but success depends on disease and nerve involvement.

\medterm{Nerve Tract} A group of nerve fibers traveling together inside the brain or spinal cord. Fibers in a tract carry signals between specific regions, so damage can cause predictable movement, sensation, or organ-control problems.

\medterm{Nervous Colon Syndrome} An older name for irritable bowel syndrome, a disorder causing recurrent abdominal pain with diarrhea, constipation, or both without a structural explanation. Symptoms may improve with individualized diet, medicines, and stress management.

\medterm{Nervous System} The brain and spinal cord (central nervous system), peripheral nerves and ganglia (peripheral nervous system), and supporting cells that detect and process information. It coordinates sensation, movement, automatic body functions, cognition, emotion, learning, and behavior.

\medterm{Nervous Tissue} Tissue made of nerve cells and supporting cells that detects changes, processes information, and sends signals through the body. It forms the brain, spinal cord, and nerves and helps control movement, sensation, and organ function.

\medterm{Nesidioblastosis} A rare diffuse or focal hyperfunctional hypertrophy and hyperplasia of pancreatic beta cells budding from the pancreatic ductal epithelium, causing endogenous persistent hyperinsulinemic hypoglycemia in neonates (congenital hyperinsulinism) and adults (post-gastric bypass hyperinsulinemic hypoglycemia).
\synonyms
Congenital Hyperinsulinism; Pancreatic Beta-Cell Hyperplasia

\medterm{NET} Abbreviation; see \textit{Neuroendocrine Tumors}.

\medterm{Netherton Syndrome} A severe autosomal recessive genodermatosis caused by mutations in the SPINK5 gene encoding the serine protease inhibitor LEKTI, characterized by the classic triad of congenital ichthyosiform erythroderma (ichthyosis linearis circumflexa), trichorrhexis invaginata ('bamboo hair'), and severe atopic diathesis.
\synonyms
SPINK5 Deficiency; Ichthyosis Linearis Circumflexa Syndrome

\medterm{Nettle Rash} An itchy rash made of raised patches that come and go, often within hours. It results from substances released in the skin and may be triggered by infection, foods, medicines, pressure, heat, or no clear cause.

\medterm{Neural} Relating to nerves, nerve cells, or the nervous system. In pregnancy, the neural tube develops into the brain and spinal cord; failure of it to close can cause birth defects such as spina bifida or anencephaly.

\medterm{Neural Crest} A temporary embryonic multipotent cellular population located along the borders of the closing neural folds that undergoes epithelial-to-mesenchymal transition and migrates extensively to form peripheral sensory ganglia, autonomic ganglia, Schwann cells, melanocytes, craniofacial bones, and adrenal chromaffin cells.

\medterm{Neuralgia} Neuralgia is sharp, burning, or electric pain along the distribution of a nerve, caused by irritation, compression, infection, injury, or other disease.

\medterm{Neuralgia Amyotrophica} An acute, idiopathic autoimmune inflammatory brachial plexopathy (Parsonage-Turner syndrome) characterized by the sudden onset of severe, unprovoked, sharp unilateral shoulder and periscapular pain, followed days to weeks later by patchy flaccid muscle weakness and rapid wasting in the distribution of the upper brachial plexus.
\synonyms
Parsonage-Turner Syndrome; Brachial Radiculitis

\medterm{Neural Tube Defects} Birth defects caused when the early structure that becomes the brain and spinal cord does not close completely during the first month of pregnancy. They include anencephaly, encephalocele, and spina bifida and vary from fatal to mild.

\medterm{Neuraminidase Inhibitors} Antiviral medicines that block a surface protein used by influenza viruses to release new virus particles from infected cells. They work best when started early and are used according to the influenza type, illness severity, and risk of complications.

\medterm{Neurapraxia} The mildest form of peripheral nerve injury in the Seddon classification, characterized by transient conduction block caused by focal demyelination or localized ischemia without structural disruption of the underlying axon or its connective tissue sheaths; complete recovery occurs spontaneously within days to weeks.

\medterm{Neurilemma} The outer layer of a Schwann cell that surrounds and supports a nerve fiber outside the brain and spinal cord.

\medterm{Neurilemmoma} Alternate spelling; see \textit{Schwannoma}.

\medterm{Neurinomatosis} Alternate name; see \textit{Schwannomatosis}.

\medterm{Neuritic Plaques} Abnormal beta-amyloid deposits surrounded by damaged nerve branches, commonly found in Alzheimer disease. They are one pathological feature of the disease and occur with other changes, including abnormal tau protein and nerve-cell loss.

\synonyms
Senile Plaques; Amyloid Plaques

\medterm{Neuritis} Inflammation or irritation of a nerve, often causing pain, altered sensation, weakness, or other changes in nerve function.

\medterm{Neuroacanthocytosis} A group of rare inherited disorders causing movement or psychiatric symptoms together with abnormal spiny red blood cells called acanthocytes. Neurologic assessment and genetic evaluation may be needed.

\medterm{Neuroblastoma} A childhood cancer arising from immature nerve cells, commonly in or near the adrenal glands or spine. It may cause an abdominal mass, pain, weakness, or other symptoms and can range from mild to aggressive.

\synonyms
Sympathicoblastoma; Neural Crest Tumor

\medterm{Neurocardiogenic Syncope} Alternate name; see \textit{Vasovagal Syncope}.

\medterm{Neurocognitive Disorder} A decline from a person’s previous level of cognitive ability, such as memory, attention, language, executive function, or judgment, that develops after normal earlier development. Mild neurocognitive disorder causes modest decline while independence is largely preserved; major neurocognitive disorder causes substantial decline that interferes with independent daily life and is often called dementia. Causes include brain disease, stroke, injury, infection, substance use, or medicines. Delirium is distinguished by its sudden onset and fluctuating attention and awareness.

\medterm{Neurocognitive Disorders} Alternate plural form; see \textit{Neurocognitive Disorder}.

\medterm{Neurocutaneous Syndrome} A disorder affecting both the nervous system and the skin, often because of an inherited genetic change. Examples include neurofibromatosis and tuberous sclerosis.

\medterm{Neurocysticercosis} Infection of the brain or spinal cord by the larval form of the pork tapeworm. It can cause seizures, headaches, confusion, or movement problems and is acquired by swallowing tapeworm eggs from contaminated food or hands.

\medterm{Neurocytoma} A usually slow-growing tumor made of nerve-cell-like cells, most often found in the fluid-filled spaces inside the brain. It may block fluid flow and cause headache, vomiting, vision changes, or balance problems; imaging and tissue testing guide care.

\medterm{Neurodegeneration} Progressive damage and loss of nerve cells, causing worsening problems with movement, memory, thinking, sensation, or other nervous-system functions.

\medterm{Neurodegeneration With Brain Iron Accumulation} A group of inherited disorders in which excess iron builds up in certain brain regions and gradually damages nerve cells. Symptoms may include abnormal movements, muscle stiffness, seizures, speech or thinking problems, and vision changes.

\medterm{Neurodegenerative Disease} A disease in which nerve cells progressively lose function or die. Examples include Alzheimer disease, Parkinson disease, and Huntington disease.

\medterm{Neurodegenerative Disorder} Another term for a condition involving progressive damage or loss of nerve cells, which can affect movement, memory, behavior, sensation, or autonomic function.

\medterm{Neurodermatitis} A chronic itchy skin condition in which repeated scratching or rubbing makes a patch thick, leathery, and irritated. Stress, dry skin, or an initial rash may start the itch–scratch cycle.
\synonyms
Lichen Simplex Chronicus; Circumscribed Neurodermatitis

\medterm{Neurodevelopmental Disorders} A group of conditions that begin during the developmental period and affect cognition, communication, behavior, motor function, learning, or adaptive skills. Examples include attention-deficit/hyperactivity disorder, autism spectrum disorder, intellectual disability, communication disorders, specific learning disorder, motor disorders, and tic disorders.

\medterm{Neuroendocrine} Relating to the connection between the nervous system and hormone system, or to cells that receive nerve signals and can release hormones. Neuroendocrine cells are found in many organs and can sometimes form tumors.

\medterm{Neuroendocrine Carcinoma} A cancer arising from cells with nerve-like and hormone-related features. It can begin in the lung or digestive system, often grows quickly, and may release hormones that cause additional symptoms.

\medterm{Neuroendocrine Carcinoma of the Skin} Alternate name; see \textit{Merkel Cell Carcinoma}.

\medterm{Neuroendocrine Tumor} A growth arising from cells that receive nerve signals and can release hormones. It may begin in organs such as the intestine, pancreas, lung, or appendix and may grow slowly or rapidly; some tumors cause hormone-related symptoms.

\medterm{Neuroendocrine Tumors} Growths arising from hormone-producing cells that also respond to nerve signals. They can occur in many organs, may release excess hormones, and range from slow-growing tumors to aggressive cancers.

\synonyms
NETs; Neuroendocrine Neoplasms; Carcinoid Tumors

\medterm{Neuroendocrine Tumors Grading} Classification of a neuroendocrine tumor by how quickly its cells divide, measured by the Ki-67 index and mitotic count. Lower grades usually grow more slowly; higher grades grow faster and may behave more aggressively. The grade helps guide treatment and prognosis.

\medterm{Neuroendocrinology} The medical field that studies how the brain and nerves control hormones and how hormones affect the nervous system. It includes conditions involving the pituitary, adrenal glands, thyroid, reproduction, growth, and body-fluid balance.

\medterm{Neurofibrillary Tangles} Twisted clumps of an abnormal protein called tau inside nerve cells. They damage the cells and are a characteristic finding in Alzheimer disease and some other disorders that cause progressive memory or thinking problems.

\synonyms
Paired Helical Filaments; Tau Tangles

\medterm{Neurofibroma} A benign peripheral nerve sheath tumor. Solitary neurofibromas are common and present as
soft, flesh-colored papules or nodules. Multiple neurofibromas are characteristic of
neurofibromatosis type 1 (NF1).

\medterm{Neurofibromatosis} A group of inherited conditions in which tumors or other changes develop in tissues containing nerves. Type 1 commonly causes café-au-lait spots, soft nerve tumors, bone changes, and learning difficulties; type 2 mainly affects hearing and balance nerves and can cause other nervous-system tumors.

\medterm{Neurofibromatosis-1} An inherited condition that can cause light-brown skin patches, freckles in the armpits or groin, soft growths along nerves, bone changes, vision problems, learning difficulties, or blood-vessel complications. Findings vary widely between people.

\medterm{Neurofibromatosis 2} An inherited condition in which tumors can grow on nerves, especially the nerves for hearing and balance. It may cause hearing loss, ringing in the ears, balance problems, cataracts, or other brain and spinal-cord tumors.

\medterm{Neurofibromatosis Type 1} An autosomal dominant inherited condition caused by a disease-causing change in the \textit{NF1} gene. It may cause light-brown skin patches, freckles in skin folds, soft growths along nerves, bone problems, vision complications, learning difficulties, or blood-vessel problems.
\synonyms
Neurofibromatosis Type I
NF1
Peripheral Neurofibromatosis
Von Recklinghausen’s Disease
Von Recklinghausen’s Neurofibromatosis
Von Recklinghausen’s Phakomatosis

\medterm{Neurofibromatosis Type 2} An autosomal dominant genetic disorder, usually caused by a disease-causing change in the \textit{NF2} gene, characterized by tumors involving nerves, especially bilateral vestibular schwannomas, with possible hearing loss, balance problems, and other nerve tumors.

\synonyms
NF2; Bilateral Acoustic Neurofibromatosis; MISME Syndrome

\medterm{Neurofibromatosis Type I} Alternate spelling; see \textit{Neurofibromatosis Type 1}.

\medterm{Neurofibromatosis Type II} Alternate spelling; see \textit{Neurofibromatosis Type 2}.

\medterm{Neurofibrosarcoma} Alternate name; see \textit{Malignant Peripheral Nerve Sheath Tumor}.

\medterm{Neurogenic} Caused by, arising from, or related to the nervous system. For example, neurogenic bladder results from nerve damage that disrupts bladder storage or emptying. The specific symptoms and treatment depend on which nerves are affected.

\medterm{Neurogenic Bladder} Bladder problems caused by damage to the nerves that control storing or passing urine. A person may have urgency, leakage, difficulty emptying, or urine retention; repeated infection and high bladder pressure can damage the kidneys.

\medterm{Neurogenic Claudication} A characteristic symptom complex of lumbar spinal canal stenosis characterized by bilateral dull aching, numbness, cramping, and heaviness in the buttocks, thighs, and calves provoked by walking or standing (lumbar extension) and relieved promptly by lumbar flexion (sitting or leaning forward over a shopping cart).
\synonyms
Pseudoclaudication; Spinal Stenosis Claudication

\medterm{Neurogenic Megacolon} Severe widening of the large intestine caused by damage to the nerves that control bowel movement. It can cause long-term constipation, abdominal swelling, and stool blockage and may lead to dangerous bowel inflammation or rupture.

\medterm{Neurogenic Orthostatic Hypotension} A drop in blood pressure after standing because nerve damage prevents the body from tightening its blood vessels and maintaining pressure. It may cause light-headedness, blurred vision, weakness, fainting, or falls.

\medterm{Neurogenic Pulmonary Edema} Rapid-onset, non-cardiogenic pulmonary edema developing acutely after severe central nervous system insults (such as massive subarachnoid hemorrhage, head trauma, or status epilepticus), provoked by an intense sympathetic surge that induces sudden pulmonary vasoconstriction and capillary leak.

\synonyms
NPE

\medterm{Neurogenic Shock} A dangerous fall in blood pressure caused by sudden loss of nerve control over blood-vessel tone, usually after a serious spinal-cord injury. The heart rate may be unusually slow and the skin warm.

\medterm{Neuroglia} Supporting cells of the nervous system that nourish, protect, insulate, and maintain nerve cells. Different types help control the nerve environment, form protective coverings, remove waste, and respond to injury or infection.

\medterm{Neurography} An MRI method that produces detailed pictures of nerves outside the brain and spinal cord. It can help assess trapped, injured, inflamed, or enlarged nerves and may help find some nerve tumors.

\medterm{Neurohypophysis} The back part of the pituitary gland that stores and releases oxytocin and vasopressin. These hormones help control childbirth, milk release, water balance, and urine concentration.

\medterm{Neurolathyrism} A toxic neurologic disorder caused by excessive consumption of Lathyrus sativus, producing progressive spastic weakness of the legs.

\medterm{Neuroleptic} An older term for an antipsychotic medicine, used to treat conditions such as schizophrenia, severe mania, or psychosis. The term remains in names such as neuroleptic malignant syndrome; these medicines may cause sleepiness, movement problems, or metabolic effects.

\medterm{Neuroleptic Agents} Antipsychotic medicines used for conditions such as schizophrenia, bipolar disorder, severe agitation, or nausea, depending on the drug. They can cause movement, metabolic, heart-rhythm, and sedation-related effects.

\synonyms
Antipsychotics

\medterm{Neuroleptic Malignant Syndrome} A rare, life-threatening reaction to some antipsychotic or other dopamine-blocking medicines. It causes very high temperature, severe muscle stiffness, confusion, sweating, unstable blood pressure, and a fast heart rate.

\medterm{Neurological} Related to the brain, spinal cord, or nerves. Neurological symptoms may include weakness, numbness, seizures, headache, dizziness, confusion, or problems with movement, speech, vision, or memory.

\medterm{Neurological Disease} Any disease affecting the brain, spinal cord, peripheral nerves, muscles, or their connections. Symptoms may include weakness, numbness, seizures, headache, movement changes, or problems with thinking.

\medterm{Neurological Examination} A clinical assessment of mental status, cranial nerves, motor function, reflexes, coordination, sensation, and gait.
It helps localize nervous-system dysfunction and determine the need for further testing.

\medterm{Neurological Rehabilitation} Treatment and training that help a person regain movement, communication, thinking, or independence after stroke, brain or spinal-cord injury, Parkinson disease, multiple sclerosis, or another nervous-system disorder.

\medterm{Neurologic Deficit} Loss or impairment of a neurologic function, such as strength, sensation, coordination, vision, language, or consciousness, caused by dysfunction of the nervous system.

\medterm{Neurologic Disorders} A group of conditions affecting the brain, spinal cord, peripheral nerves, neuromuscular junction, or muscles. Examples include stroke, epilepsy, migraine, multiple sclerosis, Parkinson disease, Alzheimer disease, peripheral neuropathy, meningitis, neuromuscular disorders, and movement disorders.

\medterm{Neurologic Localization} Using the pattern of weakness, numbness, reflex changes, coordination problems, vision, speech, or thinking changes to identify where the nervous system is affected. This helps clinicians choose likely causes and appropriate tests.

\medterm{Neurologist} A medical doctor who diagnoses and treats disorders of the brain, spinal cord, nerves, and muscles. Neurologists assess symptoms such as weakness, numbness, seizures, tremor, headaches, memory change, and movement or balance problems.

\medterm{Neurology} The medical specialty concerned with diseases of the brain, spinal cord, peripheral nerves, and muscles.

\medterm{Neurology PET} PET imaging used in selected neurologic disorders to assess brain metabolism or receptor activity, such as in some dementias, epilepsy evaluations, and movement disorders.

\medterm{Neurolysis} A procedure that frees a nerve from scar tissue or pressure, or deliberately interrupts selected nerve fibers to relieve pain. The purpose and possible effects depend on the nerve treated and the technique used.

\medterm{Neuroma} A benign growth or reactive thickening involving nerve tissue, often producing pain or altered sensation.

\medterm{Neuromodulation} Changing nerve activity with electrical, magnetic, chemical, or other signals. Treatments include spinal cord stimulation, deep brain stimulation, and some noninvasive brain-stimulation methods.

\medterm{Neuromuscular} Relating to nerves, muscles, or the place where a nerve tells a muscle to contract. Neuromuscular diseases can cause weakness, tiredness with activity, muscle wasting, cramps, or problems with breathing or swallowing.

\medterm{Neuromuscular Blockade Reversal} The pharmacological restoration of normal neuromuscular transmission following surgical paralysis, achieved by administering an acetylcholinesterase inhibitor paired with an anticholinergic (neostigmine and glycopyrrolate) or a selective encapsulating reversal agent (sugammadex for rocuronium/vecuronium).

\medterm{Neuromuscular Blocking Agents} Medicines used during anesthesia to temporarily relax skeletal muscles, making insertion of a breathing tube and some operations easier. They stop muscle movement but do not provide pain relief or unconsciousness and require breathing support and monitoring.

\medterm{Neuromuscular Disorders} A group of conditions affecting motor nerves, neuromuscular junctions, or skeletal muscles. Examples include myasthenia gravis, Lambert-Eaton syndrome, muscular dystrophy, spinal muscular atrophy, motor neuron disease, inflammatory myositis, myopathies, and peripheral neuropathies.

\medterm{Neuromuscular Hamartoma} A rare benign developmental lesion in which mature skeletal muscle and nerve fascicles are intermixed, typically involving a major peripheral nerve and causing enlargement or neuropathic symptoms.

\medterm{Neuromuscular Junction} The small connection where a movement-controlling nerve communicates with a skeletal muscle. The nerve releases a chemical messenger there, causing the muscle to contract.

\synonyms
NMJ; Motor End-Plate

\medterm{Neuromyelitis Optica} An autoimmune disorder that mainly causes inflammation of the optic nerves and spinal cord. It can cause vision loss and weakness or altered sensation and is treated with immune-directed therapy.

\medterm{Neuromyelitis Optica Spectrum Disorder} An immune-system disorder that repeatedly inflames the optic nerves, spinal cord, or parts of the brain. Attacks may cause vision loss, limb weakness, numbness, vomiting, or bladder problems and can lead to lasting disability without treatment.

\synonyms
Neuromyelitis Optica
NMO
NMOSD
Devic Disease
Devic's Disease

\medterm{Neuromyotonia} A rare disorder in which nerves repeatedly stimulate muscles, causing ongoing stiffness, twitching, cramps, delayed muscle relaxation, and sometimes excessive sweating. It may be inherited or caused by an immune reaction. An electrical nerve and muscle test helps diagnosis; treatment aims to calm nerve activity and address any underlying cause.

\synonyms
Isaacs Syndrome; Continuous Muscle Fiber Activity Syndrome

\medterm{Neuron} A nerve cell that receives, processes, and sends messages using electrical signals and chemical messengers. It has a cell body, branches that receive information, and usually a long fiber that carries messages to other cells.

\medterm{Neuronal Ceroid Lipofuscinoses} A group of inherited disorders in which waste material builds up inside nerve cells. They usually cause worsening seizures, vision loss, movement problems, and difficulty thinking or learning, with symptoms and age of onset varying by type.

\medterm{Neuronal Migration Disorder} A developmental brain disorder in which nerve cells do not move to their normal positions before birth. It can cause abnormal brain structure, seizures, developmental delay, or learning problems.

\medterm{Neurone} The British spelling of neuron, an excitable nerve cell that receives and transmits signals.

\medterm{Neuronopathy} A disorder caused mainly by damage to nerve-cell bodies, often in sensory or motor nerve groups. It may cause numbness, imbalance, weakness, or loss of reflexes.

\medterm{Neurons} Specialized nerve cells that receive, process, and transmit electrical or chemical signals. They communicate with one another and with muscles or glands.

\medterm{Neuron-Specific Enolase} A protein that can be released by nerve cells and some neuroendocrine tumors. A raised blood level may support evaluation or monitoring in selected situations but is not specific enough to diagnose a disease alone.

\medterm{Neuron-Specific Enolase Tumor Marker} Neuron-specific enolase, a protein that may be increased in some neuroendocrine cancers or after nerve-cell injury. It can support monitoring in selected situations but is not specific enough to diagnose cancer alone.

\medterm{Neuro-Ophthalmology} The medical specialty that evaluates vision problems caused by the brain, optic nerves, or eye-movement pathways. It assesses symptoms such as unexplained vision loss, double vision, abnormal pupils, or visual-field changes.

\medterm{Neuropathic} Caused by disease or dysfunction of nerves; neuropathic skin symptoms may include burning, electric-shock pain, altered sensation, numbness, or itch.

\medterm{Neuropathic Pain} Pain caused by disease or injury of the nerves that carry sensation, the spinal cord, or the brain. It is often burning, shooting, electric, or accompanied by numbness and may follow diabetes, shingles, injury, or nerve compression.

\medterm{Neuropathic Ulcer} A wound caused or worsened by loss of protective feeling, most often in the foot. A person may not notice pressure or injury, allowing the wound to deepen; diabetes, poor circulation, and infection can delay healing and increase the risk of amputation.

\medterm{Neuropathy} Disease or damage affecting nerves outside the brain and spinal cord. It may cause burning pain, numbness, tingling, weakness, balance problems, or altered sweating and can result from diabetes, medicines, toxins, poor nutrition, compression, or inherited disease.

\medterm{Neuropeptide} A peptide released by neurons that acts as a signaling molecule in the nervous or endocrine system.

\medterm{Neuropeptides} Short chains of amino acids released by nerve cells as chemical signals. They can influence pain, appetite, stress, mood, sleep, blood pressure, inflammation, hormone release, and communication between nerve cells.

\medterm{Neuropharmacology} The study of how medicines and other substances affect the brain, spinal cord, and nerves. It examines their actions, benefits, side effects, interactions, and uses in treating neurological or psychiatric conditions.

\medterm{Neurophysiotherapy} Physical therapy focused on movement and function problems caused by disorders of the brain, spinal cord, or nerves. Treatment may address strength, balance, walking, coordination, and daily activities.

\medterm{Neuroplasticity} The brain’s ability to change its connections and activity in response to learning, experience, or injury. This process supports recovery and skill development, especially with repeated practice and rehabilitation, but recovery varies with the injury and the person.

\medterm{Neuroplegic} Causing paralysis or severe loss of nerve-controlled movement. The effect may result from a neurologic disease, spinal injury, toxin, or medicine and can impair breathing or other vital functions when major nerve pathways are affected.

\medterm{Neuropsychiatry} The clinical field linking neurologic and psychiatric disorders, including cognitive and behavioral consequences of brain disease.

\medterm{Neuroradiology} The use of imaging to diagnose disorders of the brain, spine, head and neck, and peripheral nerves.

\medterm{Neuroretinitis} Inflammation of the optic nerve and nearby retina that can cause blurred or reduced vision, blind spots, and pain with eye movement. Infections and immune disorders are possible causes and require specialist evaluation.

\medterm{Neurosarcoidosis} Sarcoidosis affecting the brain, spinal cord, coverings around them, or peripheral nerves. It may cause headaches, facial weakness, vision problems, seizures, weakness, numbness, or hormone disturbances and usually needs specialist treatment.

\medterm{Neurosciences} Neurosciences are the scientific disciplines studying the nervous system, including its structure, function, development, disease, and treatment.

\medterm{Neurosurgeon} A surgeon trained to diagnose and operate on disorders of the brain, spinal cord, spine, nerves, and nearby supporting structures. Neurosurgeons may treat tumors, injuries, bleeding, pressure, pain, or blood-vessel problems.

\medterm{Neurosurgery} Surgery for disorders of the brain, spinal cord, spine, peripheral nerves, or supporting structures. Procedures may remove tumors, relieve pressure, repair injuries, stop bleeding, treat blood vessels, or improve severe pain.

\medterm{Neurosyphilis} Neurosyphilis is infection of the nervous system by Treponema pallidum and can cause meningitis, stroke, sensory loss, cognitive change, or tabes dorsalis.

\medterm{Neurotoxicity} Injury or dysfunction of the nervous system caused by a medicine, drug, toxin, heavy metal, infection, or other harmful exposure.

\medterm{Neurotransmission} The process by which nerve cells communicate with one another or with muscles and glands using electrical signals and chemical messengers. It involves release of a neurotransmitter, movement across a synapse, receptor binding, and termination or removal of the signal.

\medterm{Neurotransmitter} A chemical messenger released by a nerve cell. It crosses a tiny gap to affect another nerve cell or a muscle or gland cell, either increasing or reducing its activity.

\medterm{Neurotransmitter Receptors} Proteins on or inside cells that receive chemical messages released by nerve cells. When a message attaches to a receptor, it can increase or decrease activity, affecting movement, mood, pain, heart rate, or other body functions.

\medterm{Neurotrophic} Describing a substance or effect that supports the growth, survival, connection, or repair of nerve cells. Neurotrophic signals help nerves develop and maintain function and are studied in nerve injury and degenerative disease.

\medterm{Neurotrophic Keratopathy} Corneal damage caused by reduced sensation from injury or disease of the trigeminal nerve. It can lead to poor healing, recurrent epithelial defects, ulceration, and vision loss, sometimes with little pain.

\medterm{Neurotropic} Having a preference for nervous tissue or mainly affecting the brain, spinal cord, or nerves. The term is often used for infections, medicines, or other substances that act in these tissues.

\medterm{Neurovascular Examination} A check of movement, feeling, pulses, skin color, temperature, and blood return in an injured arm or leg. Repeated checks can detect worsening nerve damage, reduced blood flow, or dangerous swelling and pressure.

\medterm{Neutral Alignment} A posture in which the body is arranged in a straight, balanced line while preserving normal spinal curves.

\medterm{Neutral Endopeptidase} A membrane-bound zinc metallopeptidase expressed in vascular endothelial and renal proximal tubular cells that degrades circulating natriuretic peptides (ANP, BNP, CNP), bradykinin, and substance P; pharmacologically inhibited by sacubitril in heart failure therapy.

\synonyms
Neprilysin; NEP; CD10

\medterm{Neutralization} Removal or reduction of an effect by counteraction, such as an acid-base reaction or antibody-antigen binding.

\medterm{Neutralizing Antibodies} Antibodies that bind a pathogen or its toxin and block its ability to enter cells, attach to tissues, or cause its usual biologic effects.

\medterm{Neutropenia} An abnormally low number of neutrophils, white blood cells that help fight bacterial and fungal infections. The lower the count, the greater the infection risk; fever during severe neutropenia can signal a serious infection.

\synonyms
Neutrocytopenia; Low Absolute Neutrophil Count

\medterm{Neutropenic Enterocolitis} A life-threatening, necrotizing transmural inflammation of the cecum and ascending colon occurring in severely neutropenic oncology patients receiving intensive myelosuppressive chemotherapy, presenting with fever, right lower quadrant abdominal pain, and bowel wall thickening on CT.

\synonyms
Typhlitis

\medterm{Neutropenic Fever} A medical emergency in oncology patients defined as a single oral temperature of 38.3 degrees Celsius (101 F) or sustained temperature of 38.0 degrees Celsius for one hour in the setting of an absolute neutrophil count (ANC) below 500 cells/microliter, requiring immediate empiric broad-spectrum antipseudomonal antibiotic administration.

\synonyms
Febrile Neutropenia

\medterm{Neutrophil} The most common infection-fighting white blood cell in the bloodstream. It quickly enters injured tissue, surrounds and destroys many bacteria and fungi, and helps start inflammation; unusually low levels increase infection risk.

\medterm{Neutrophil Extracellular Traps} Web-like material released by neutrophils, made mainly of DNA and antimicrobial proteins. These traps can help contain germs, but excessive formation may injure tissues and contribute to inflammation, clotting, or autoimmune disease.

\medterm{Neutrophil Extracellular Traps} Webs of extracellular decondensed chromatin fibers decorated with antimicrobial granule proteins (myeloperoxidase, elastase, histones) extruded by activated neutrophils during a specialized form of cell death (NETosis) to trap and kill extracellular bacterial pathogens; implicated in thrombosis and autoimmune disease.

\synonyms
NETs

\medterm{Neutrophilic Dermatosis} A group of inflammatory skin diseases in which infection-fighting white blood cells collect in the skin without causing a typical skin infection. Examples include Sweet syndrome and pyoderma gangrenosum; the rash may be painful, swollen, or associated with fever.

\medterm{Neutrophilic Eccrine Hidradenitis} An uncommon skin inflammation in which infection-fighting white cells collect around sweat glands. It often occurs during chemotherapy and causes painful red or purple bumps or patches, usually on the trunk or limbs.

\medterm{Neutrophil Recruitment} The movement of neutrophils from the bloodstream into injured or infected tissue. They first slow and attach to the blood-vessel lining, then pass through it and follow chemical signals to the affected area.

\medterm{Neutrophils} The most abundant type of white blood cell. Neutrophils rapidly migrate to sites of infection or injury and ingest microorganisms and cellular debris.

\medterm{Nevoid Basal Cell Carcinoma Syndrome} Gorlin syndrome, an inherited disorder increasing the risk of basal-cell cancers, jaw cysts, and other developmental abnormalities.

\medterm{Nevus} A usually harmless, localized change or growth in the skin or mucous membrane, often called a mole. It may contain pigment-producing cells or involve blood vessels or other skin structures; a changing lesion needs assessment.

\medterm{Nevus Anemicus} A congenital vascular anomaly presenting as a localized, well-circumscribed, pale, hypopigmented macule or patch that remains pale after rubbing or temperature changes due to localized hypersensitivity of cutaneous alpha-adrenergic receptors producing constant arteriolar vasoconstriction.
\synonyms
Anemic Nevus; Pharmacologic Vasoconstrictive Nevus

\medterm{Nevus Comedonicus} A congenital hamartoma of the pilosebaceous unit presenting as a localised cluster of
dilated follicular ostia filled with keratinous plugs, resembling closely grouped open
comedones. It typically appears at birth or in early childhood on the face, neck, or
trunk, and may be associated with other developmental anomalies.

\medterm{Nevus Comedonicus Syndrome} A rare disorder in which multiple comedone-like skin lesions occur with abnormalities of the eyes, bones, nervous system, or other organs. It is usually present from birth or early childhood.

\medterm{Nevus of Ito} A benign dermal melanocytic hamartoma characterized by mottled, patchy slate-gray, blue-brown, or hyperpigmented macules distributed along the sensory branches of the posterior supraclavicular and lateral cutaneous brachial nerves (involving the supraclavicular, acromial, and scapular regions).
\synonyms
Omohyoid Nevus; Dermal Melanocytosis of Shoulder

\medterm{Nevus of Ota} A congenital or acquired dermal melanocytosis presenting as mottled blue-gray or brownish-black pigmentation in the distribution of the ophthalmic (V1) and maxillary (V2) branches of the trigeminal nerve, commonly involving the ipsilateral sclera, conjunctiva, and facial skin.
\synonyms
Oculodermal Melanocytosis; Nevus Fusco-Caeruleus Ophthalmomaxillaris

\medterm{Nevus Sebaceus} A congenital or early-onset hamartoma of sebaceous and other adnexal structures, usually a yellow-orange hairless plaque on the scalp or face; secondary tumors may rarely arise within it.

\medterm{Newborn} A baby during the first 28 days after birth, when special newborn health monitoring and care are important.

\medterm{Newborn Circumcision} Surgical removal of the foreskin from the penis, typically performed within the first
few days of life. Medical benefits include reduced risk of urinary tract infections,
sexually transmitted infections, and penile cancer. Complications include bleeding,
infection, and meatal stenosis.

\medterm{Newborn Head Molding} Temporary shaping or overlapping of a newborn’s skull bones during passage through the birth canal; it usually resolves as the skull returns to its normal contour.

\medterm{Newborn Hearing Screen} A quick test, usually before discharge from the birth facility, that checks whether a newborn responds normally to sound. A result needing repeat testing does not prove hearing loss; timely follow-up is important when the screen is not passed.

\medterm{Newborn Infant Development} The normal physical, neuromotor, sensory, and social milestone progression of a healthy full-term infant during the first month of life (neonatal period), including primitive reflexes (Moro, rooting, suckling, palmar grasp), visual tracking, auditory responsiveness, and weight recovery.

\synonyms
Neonatal Development; Infant Milestone Progression; Newborn Neuromotor Development

\medterm{Newborn Jaundice} Newborn jaundice is yellow discoloration from elevated bilirubin; physiologic patterns are common, but early, severe, rapidly rising, or conjugated jaundice requires evaluation for serious disease.

\medterm{Newborn Metabolic Screen} A dried blood spot test (heel stick) collected 24-48 hours after birth to screen for
congenital disorders before symptoms develop.

\medterm{Newborn Screening} Testing shortly after birth for selected treatable genetic, metabolic, hormone, hearing, and heart conditions before symptoms appear. A small blood sample is usually taken from the heel; the conditions screened and timing vary by country or region.

\medterm{Newborn Screening Panel} A standardized public health laboratory testing program performed on dried blood spots collected via neonatal heel prick within 24 to 48 hours of life, screening for early, treatable inborn errors of metabolism, hemoglobinopathies, and endocrine disorders (such as phenylketonuria, congenital hypothyroidism, cystic fibrosis, and sickle cell disease).

\medterm{Newborn Screening Tests} Tests performed shortly after birth to detect selected treatable genetic, endocrine, metabolic, hematologic, hearing, or critical congenital heart conditions before symptoms develop.

\medterm{New Daily Persistent Headache} A headache that begins on a clearly remembered day and becomes continuous within 24 hours, persisting for
at least three months. It may resemble migraine or tension-type headache; a new persistent headache requires
medical evaluation to exclude secondary causes such as infection, pressure disorders, or vascular disease.

\medterm{New York Heart Association Functional Classification} A standardized clinical grading scale assessing the severity of functional symptomatic limitation in patients with heart failure based on physical exertion, ranging from Class I (no symptoms with ordinary physical activity) to Class IV (symptoms present at rest).

\synonyms
NYHA Functional Classification; NYHA Class

\medterm{Next-Generation Sequencing} A laboratory method that reads many DNA sections at the same time. It can examine selected genes, nearly all protein-coding genes, the whole genome, or tumor DNA; results may miss some changes and require expert interpretation.

\medterm{Nexus} A connection or central point linking structures, processes, or ideas. In medical writing, its meaning depends on the context and may describe an anatomic connection, a group of related causes, or a link between two findings.

\medterm{NF1} Abbreviation; see \textit{Neurofibromatosis Type 1}.

\medterm{NF2} Abbreviation; see \textit{Neurofibromatosis Type 2}.

\medterm{Niacin} Vitamin B3, needed to help cells release energy and perform many chemical reactions; it forms the coenzymes NAD+ and NADP+ that carry electrons in metabolism. Severe deficiency causes pellagra, with skin changes, digestive problems, confusion, weakness, and potentially life-threatening complications.

\synonyms
Vitamin B3; Nicotinic Acid

\medterm{Niacinamide} The amide form of vitamin B3, used as a nutritional supplement and in selected dermatologic preparations. Unlike niacin, it usually does not cause flushing at nutritional doses.

\medterm{Niacin Deficiency} A lack of vitamin B3, usually from a diet poor in niacin and
the amino acid tryptophan. When severe it causes pellagra, with dermatitis, diarrhea, and
neuropsychiatric disturbance, and it may complicate alcohol use disorder or conditions
that impair absorption. Niacin replacement and treatment of the underlying cause correct
the deficiency.

\medterm{Niacin For Cholesterol} Niacin lowers triglycerides and can raise HDL cholesterol, but its routine cardiovascular benefit is limited when added to effective statin therapy; flushing, liver injury, hyperglycaemia, and hyperuricaemia limit high-dose use.

\medterm{Niacin Vitamin B3} Alternate name; see \textit{Niacin}.

\medterm{Nialamide} An older monoamine-oxidase-inhibitor antidepressant that is no longer commonly used.

\medterm{Nicardipine} A calcium-channel-blocking medicine that relaxes blood vessels and lowers blood pressure. It is used for selected urgent or ongoing blood-pressure problems and may cause headache, flushing, dizziness, fast heartbeat, or swelling.

\medterm{Nickel Allergy} An allergic skin reaction caused by contact with nickel-containing materials.

\synonyms
Nickel Hypersensitivity; Nickel Contact Dermatitis

\medterm{Nickel-Chromium} A metal mixture used in some dental crowns and medical devices. It is strong and less costly than some alternatives, but people allergic to nickel may develop skin or mouth reactions after contact.

\medterm{Niclosamide} An antiparasitic medicine used in some countries to treat selected intestinal tapeworm infections. It acts mainly in the gut and does not reliably treat larval tapeworm infections in tissues, such as neurocysticercosis.

\medterm{Nicolau Syndrome} A rare, acute iatrogenic cutaneous adverse reaction characterized by severe localized pain, livedoid erythema, and rapid ischemic skin necrosis (embolia cutis medicamentosa) occurring at the site of intramuscular or intra-articular drug injections, caused by accidental intra-arterial or peri-arterial drug deposition and microvascular thrombosis.
\synonyms
Embolia Cutis Medicamentosa; Livedo-like Dermatitis

\medterm{Nicotinamide} Alternate name; see \textit{Niacinamide}.

\medterm{Nicotinamide Adenine Dinucleotide} A molecule present in all living cells that helps transfer energy from food and supports enzymes involved in cell repair and communication. It exists in forms commonly written as NAD+ and NADH.

\medterm{Nicotine} An addictive stimulant found in tobacco and some vaping products that acts on neuronal nicotinic receptors.

\medterm{Nicotine Addiction} Alternate name; see \textit{Nicotine Dependence}.

\medterm{Nicotine And Tobacco} Nicotine poisoning can follow swallowing tobacco or liquid nicotine, breathing concentrated vapor, or skin contact. It may cause nausea, vomiting, sweating, dizziness, fast heartbeat, seizures, or breathing failure.

\medterm{Nicotine Dependence} A pattern of compulsive nicotine use with cravings, loss of control, tolerance, and withdrawal symptoms when use is reduced or stopped.

\synonyms
Tobacco Use Disorder; Nicotine Addiction

\medterm{Nicotine Poisoning} Nicotine poisoning causes early nausea, vomiting, salivation, dizziness, tachycardia, and agitation followed by weakness, bradycardia, seizures, or respiratory failure in severe exposure.

\medterm{Nicotine Replacement Therapy} Nicotine medicines help reduce withdrawal when stopping tobacco. Using too much or combining products improperly can cause nausea, vomiting, dizziness, sweating, fast heartbeat, confusion, or seizures; keep them away from children.

\medterm{Nicotinic} Relating to nicotine or to nicotinic acetylcholine receptors. These receptors respond to the chemical messenger acetylcholine and help control nerve signals, skeletal-muscle movement, attention, and some automatic body functions.

\medterm{Nicotinic Acetylcholine Receptor} A pentameric ligand-gated non-selective cation channel responsive to acetylcholine and nicotine, functioning to generate rapid excitatory postsynaptic potentials at the neuromuscular junction (muscle type) and throughout autonomic ganglia and the central nervous system (neuronal type).

\synonyms
nAChR

\medterm{Nicotinic Acid} Alternate name; see \textit{Niacin}.

\medterm{NICU Consultants And Support Staff} NICU consultants and support staff are specialists and professionals who help care for critically ill or premature newborns, including physicians, nurses, therapists, and social workers.

\medterm{NICU Staff} NICU staff are the healthcare professionals who monitor and treat newborns in a neonatal intensive care unit and support their families.

\medterm{Nidation} Implantation of a fertilized egg into the lining of the uterus early in pregnancy. The term is largely synonymous with implantation.

\medterm{Nidus} A starting point or central focus from which a process develops. In medicine, it may mean the core of a kidney stone, the focus of an infection, or the central part of an abnormal blood-vessel network.

\medterm{Niemann-Pick Disease} A group of rare inherited lysosomal storage disorders in which lipids accumulate in cells. The group includes acid sphingomyelinase deficiency, formerly called Niemann-Pick disease types A and B, and Niemann-Pick disease type C, caused by impaired intracellular cholesterol trafficking. Features vary by type and may include enlargement of the liver or spleen, lung disease, difficulty coordinating movement, and progressive neurologic impairment.

\medterm{Nifedipine} A medicine that relaxes blood vessels to lower blood pressure and relieve some types of chest pain. It may cause flushing, headache, dizziness, ankle swelling, or a fast heartbeat; the dose depends on the condition and formulation.

\medterm{Night Blindness} Difficulty seeing in dim light or darkness. Causes include vitamin-A deficiency, inherited retinal disorders, cataracts, glaucoma, and some medicines. New or worsening night blindness needs an eye examination because some causes can be treated.

\medterm{Night-Blindness} Difficulty seeing in dim light, called nyctalopia; causes include vitamin A deficiency, retinal disease, and inherited disorders.

\medterm{Night Eating Syndrome} A recurring pattern of eating much of the day’s food in the evening or after waking at night, often with little morning appetite and sleep problems. It is a related eating pattern that may occur with depression, obesity, or other eating disorders and is not automatically a formal eating-disorder diagnosis.

\medterm{Night Leg Cramps} Sudden, painful involuntary contractions of leg muscles occurring during sleep or rest, often affecting the calf or foot.

\medterm{Nightmare Disorder} Repeated frightening dreams that cause distress, wake a person, or interfere with sleep and daytime activities. It may be linked with stress, trauma, medicines, mental health conditions, or other sleep disorders.

\medterm{Nightmares} Vivid frightening dreams that cause distress or wake a person, sometimes linked to stress, illness, or medicines.

\medterm{Night Sweats} Episodes of excessive sweating during sleep that soak clothing or bedding. They may occur with infection,
hormonal changes, cancer, medication effects, or other systemic illness; persistent unexplained
night sweats warrant assessment.

\synonyms
Sleep Hyperhidrosis; Nocturnal Diaphoresis

\medterm{Night-Sweats} Episodes of excessive sweating during sleep; causes include infection, malignancy, hormonal change, medicines, and environmental heat.

\medterm{Night Sweats in Men} Alternate name; see \textit{Night Sweats}.

\medterm{Night Terrors} Episodes of abrupt partial awakening from deep non-REM sleep with intense fear, screaming, autonomic arousal, and limited recall afterward, occurring most often in children.

\medterm{Night Terrors In Children} Night terrors are episodes of sudden fear, screaming, autonomic arousal, and apparent confusion during non-REM sleep; children usually remain difficult to awaken and have little memory afterward.

\medterm{NIH Stroke Scale} The National Institutes of Health Stroke Scale is a standardized, validated
clinical assessment tool used to quantify the severity of acute ischemic stroke.

\medterm{Nijmegen Breakage Syndrome} A rare inherited disorder that affects DNA repair and the immune system. It may cause a small head, short stature, developmental or learning difficulties, frequent infections, and a high risk of lymphoma or other cancers. People with it are unusually sensitive to radiation, so genetic and specialist care is important.

\synonyms
NBS; Ataxia-Telangiectasia Variant 1

\medterm{Nikethamide} An old medicine once used to stimulate breathing. It is no longer a routine treatment because its benefit is limited and excessive doses can cause agitation, shaking, seizures, or other dangerous effects.

\medterm{Nikolsky Sign} Peeling or separation of the top skin layer when nearby skin is gently rubbed. It can occur in serious blistering disorders, but the sign is not specific and must be interpreted with the rash and other findings.

\medterm{Nikolsky's Sign} Shearing or peeling of the superficial epidermis when apparently normal skin beside a blister is rubbed; it may occur in pemphigus and other blistering disorders.

\medterm{Nipah Virus} A zoonotic virus in the genus \textit{Henipavirus}, with fruit bats as its natural host. It can spread from infected bats or other animals to people, through contaminated food such as raw date-palm sap, and between people through close contact. Infection ranges from no symptoms to fever, headache, vomiting, and respiratory illness such as cough or difficulty breathing. Severe disease may cause inflammation of the brain (encephalitis), confusion, seizures, coma, or death. There is currently no licensed vaccine or specific approved treatment for Nipah virus disease; management relies on supportive care.

\medterm{Nipah Virus Infection} A zoonotic illness caused by Nipah virus. Fruit bats are the natural host, and infection can pass to people through infected animals, contaminated food such as raw date-palm sap, or close contact with an infected person. Some infections cause no symptoms; illness may otherwise cause fever, headache, muscle pain, vomiting, cough, or difficulty breathing and can progress to brain inflammation (encephalitis), confusion, seizures, coma, or death. Early intensive supportive care may improve survival, but there is currently no licensed vaccine or specific approved treatment.

\medterm{Nipple} The projection at the centre of the breast through which milk is released from the mammary ducts; it contains smooth muscle and sensory nerve endings.

\medterm{Nipple Discharge} Fluid leaking from one or both nipples when not pregnant or breastfeeding. Most nipple discharge is harmless, coming from multiple milk ducts in both breasts only when squeezed, and is usually caused by normal hormone changes or medications. Discharge that leaks on its own from a single duct in one breast, especially if clear or bloody, is less common. The most frequent cause is a small, harmless growth inside a milk duct (a papilloma), while breast cancer accounts for about 5\% to 15\% of these cases. Doctors identify the cause using a physical examination, imaging such as ultrasound or mammography, and duct tests.

\medterm{Pathologic Nipple Discharge} Nipple fluid that occurs without squeezing, comes from one breast or one duct, or is bloody or clear. It may have a benign cause, but it needs clinical assessment and often imaging to decide whether further testing is needed.

\medterm{Nipple Inversion Or Retraction} Inward turning or pulling-in of the nipple. It may be longstanding and benign, but new one-sided inversion or retraction can occur with duct disease or breast cancer.

\medterm{Nipple Retraction} Alternate name; see \textit{Nipple Inversion Or Retraction}.

\medterm{Niridazole} An older medicine once used against schistosomiasis, a parasitic infection. Its use is now limited because of serious possible side effects, including seizures, liver injury, and damage to blood-cell production.

\medterm{Nissen Fundoplication} A surgical antireflux procedure in which the gastric fundus is completely wrapped 360 degrees around the lower esophagus to reinforce the lower esophageal sphincter (LES) and repair a hiatal hernia, used for refractory gastroesophageal reflux disease (GERD).
\synonyms
360-Degree Complete Fundoplication; Total Antireflux Fundoplication

\medterm{Nit} The egg of a louse, usually attached firmly to a hair shaft or clothing fiber. Nits may be confused with dandruff, but they do not brush away easily. Finding live lice confirms an active infestation; treatment and fine combing remove lice and eggs.

\medterm{Nitabuch Layer} A normal, dense, fibrinoid layer formed at the junction between the invading syncytiotrophoblast and the decidua basalis (stratum compactum) in normal pregnancy, which acts as a physiological anatomical boundary limiting deep placental invasion (congenitally deficient in placenta accreta).
\synonyms
Stria of Nitabuch; Decidual Fibrinoid Boundary

\medterm{Nitrates} Medicines that widen blood vessels by releasing nitric oxide. They can relieve angina and lower the heart's workload but may cause headache, flushing, dizziness, or dangerously low blood pressure with certain medicines.

\medterm{Nitrate Tolerance} Reduced effect of nitrate medicines, such as nitroglycerin, after continuous use. It may involve changes in blood-vessel signaling and the body’s counter-response to the medicine.

\synonyms
Nitroglycerin Tachyphylaxis; Nitrate Resistance

\medterm{Nitrazepam} A benzodiazepine medicine that slows brain activity and can cause sleepiness, calmness, and seizure prevention. It can cause dependence, falls, memory problems, and dangerously slow breathing, especially with alcohol or opioids.

\medterm{Nitric Acid} A strong corrosive chemical that can burn skin, eyes, mouth, throat, and digestive organs. Breathing its fumes can injure the lungs, and swallowed or inhaled exposure can cause severe poisoning.

\medterm{Nitric Acid Poisoning} Exposure to nitric acid can severely burn the skin, eyes, mouth, throat, and stomach and may damage the lungs if fumes are inhaled.

\medterm{Nitric Oxide} A short-lived signaling molecule made by many cells. It relaxes blood-vessel muscle, helping regulate blood flow and pressure, and also participates in nerve signaling and immune defense. Too much or too little production may contribute to disease.

\medterm{Nitric Oxide Synthase} An enzyme family that catalyzes the synthesis of the gaseous signaling molecule nitric oxide (NO) and L-citrulline from L-arginine, oxygen, and NADPH, comprising constitutive calcium-dependent isoforms (endothelial eNOS, neuronal nNOS) and a calcium-independent inducible isoform (iNOS).

\synonyms
NOS

\medterm{Nitrofurantoin} An antibacterial medicine used mainly for bladder infections caused by susceptible bacteria. It concentrates in urine but is not suitable for many kidney infections; nausea, lung reactions, liver injury, or nerve problems are uncommon but important risks.

\medterm{Nitrogen} An essential element in amino acids, proteins, nucleic acids, and other biological molecules; atmospheric nitrogen is largely inert to humans.

\medterm{Nitrogen Balance} The difference between dietary nitrogen intake, estimated mainly from protein intake,
and nitrogen loss, primarily in urine. Positive balance supports growth or tissue
repair; negative balance occurs with inadequate intake, inflammation, trauma, or
catabolic disease.

\medterm{Nitrogenous Base} A nitrogen-containing component of nucleic acids; adenine, guanine, cytosine, thymine, and uracil pair in DNA or RNA to encode genetic information.

\medterm{Nitroglycerin Overdose} Too much nitroglycerin can cause severe headache, flushing, dizziness, fainting, low blood pressure, fast heartbeat, or breathing difficulty.

\medterm{Nitromors} A trade name for paint or varnish stripper formulations; exposure can cause chemical burns or toxic inhalation.

\medterm{Nitrosoureas} A group of chemotherapy medicines that can cross the blood-brain barrier and damage DNA; examples include carmustine and lomustine.

\medterm{Nitrous Oxide} An inhaled anesthetic gas with analgesic and anesthetic-sparing properties. It has rapid
onset and offset but relatively low potency, expands closed gas spaces, and may
contribute to diffusion hypoxia or postoperative nausea.

\medterm{Nitrous Oxide For Labor} A 50:50 mixture of nitrous oxide and oxygen (Entonox) inhaled by the mother during labor
for analgesia. Self-administered via a mask or mouthpiece; the woman controls her own
dosing, taking a breath 30 seconds before a contraction.

\medterm{Nizatidine} A medicine that reduces stomach-acid production by blocking histamine signals in stomach cells. It may relieve ulcers or reflux symptoms, although availability varies and treatment should consider kidney function and medicine interactions.

\medterm{NMR} NMR, or nuclear magnetic resonance, measures how atomic nuclei respond to a magnetic field and is used in chemical analysis and imaging.

\medterm{Nocardia Asteroides Infection} An infection caused by Nocardia asteroides, a soil-associated bacterium. It most often affects the lungs, skin, or brain and is more likely when immunity is weakened; treatment requires species identification and prolonged antibiotics.

\medterm{Nocardia Infection} An infection caused by Nocardia bacteria, usually affecting the lungs, skin, or brain, especially in people with weakened immunity.

\medterm{Nocardiosis} An infection caused by Nocardia bacteria found in soil and water. It most often affects the lungs, skin, or brain, especially in people with weakened immunity, and may cause cough, fever, skin sores, confusion, or seizures.

\medterm{Nocebo} A harmful or unpleasant symptom or treatment effect caused or amplified by negative expectations rather than by the treatment’s specific pharmacologic action.

\medterm{Nociception} The nervous system's detection and processing of potentially harmful heat, pressure, or chemicals. It can lead to pain, but pain also depends on how the brain interprets these signals.

\medterm{Nociceptor} A sensory nerve ending that detects potentially damaging pressure, heat, cold, or chemicals. It sends warning signals through nerves to the spinal cord and brain, where they can be experienced as pain.

\medterm{Nociceptors} Sensory nerve endings that detect potentially damaging pressure, heat, cold, or chemicals. They send warning signals through the nervous system, although pain also depends on how the brain interprets those signals.

\medterm{Nocturia} Waking during the night one or more times to urinate and then returning to sleep. It may result from evening fluids, medicines, bladder or prostate problems, diabetes, heart disease, or sleep disorders.

\medterm{Nocturnal} Occurring at night or during sleep. Examples include bedwetting, nighttime cough, breathing pauses during sleep, nighttime low oxygen, or breathlessness that wakes a person. The cause depends on the symptom and requires assessment when persistent or severe.

\medterm{Nocturnal Emission} Involuntary ejaculation during sleep, sometimes called a wet dream. It is a common and normal event during male puberty and adolescence.

\medterm{Nocturnal Enuresis} Repeated involuntary urination during sleep in a child who is old enough to have developed nighttime bladder control. It may run in families; pain, daytime wetting, constipation, or sudden onset needs evaluation.
\synonyms
Nocturnal Incontinence; Sleep Enuresis

\medterm{Nocturnal Enuresis Entry} Alternate name; see \textit{Bed-Wetting}.

\medterm{Nocturnal Hypoventilation} A reduction in minute ventilation occurring specifically during sleep, resulting in nighttime hypercapnia and hypoxemia, commonly encountered in neuromuscular diseases (such as ALS or muscular dystrophies), severe kyphoscoliosis, and obesity hypoventilation syndrome.

\medterm{Nocturnal Incontinence} Alternate name; see \textit{Incontinence}.

\medterm{Nocturnal Lagophthalmos} Incomplete closure of the eyelids during sleep, which may cause dryness, irritation, and exposure injury to the cornea.

\medterm{Nocturnal Polyuria} Production of an unusually large share of the day’s urine during nighttime sleep. It can contribute to repeated nighttime urination and may be related to fluid shifts, sleep disorders, medicines, or other medical conditions.

\medterm{Nodal Rhythm} A heartbeat rhythm that starts near the atrioventricular node instead of the heart’s usual pacemaker. It may occur when the normal pacemaker is slow or blocked and can cause a slow pulse, dizziness, or no symptoms.

\medterm{Node} A localized swelling, structure, or point of connection; examples include lymph nodes and electrical conduction nodes.

\medterm{Node of Ranvier} Periodic, unmyelinated gaps or interruptions occurring at regular intervals along the length of a myelinated nerve axon where high concentrations of voltage-gated sodium channels enable rapid saltatory nerve impulse conduction.
\synonyms
Ranvier Node; Myelin Sheath Gap

\medterm{Nodular} Made up of or containing one or more rounded lumps called nodules. The word may describe a skin lesion, thyroid swelling, lung finding, or other tissue. Its importance depends on the nodule’s size, appearance, location, growth, and associated symptoms.

\medterm{Nodular Fasciitis} A rapidly growing benign self-limited fibroblastic and myofibroblastic proliferation, usually presenting as a painful or tender superficial mass. It can mimic sarcoma clinically and microscopically.

\medterm{Nodular Lymphoid Hyperplasia} Enlargement of small areas of immune tissue, most often in the small intestine. It is usually harmless but may be associated with reduced antibody production or intestinal infection; its significance depends on symptoms and the test findings.

\medterm{Nodular Melanoma} A type of skin cancer that grows downward early and often appears as a rapidly enlarging firm bump. It may be black, brown, pink, or skin-colored, so any quickly changing or bleeding lump needs prompt assessment.

\medterm{Nodule} A small, rounded area of abnormal tissue that may be solid or fluid-filled and may occur in the skin, lung, thyroid, or another organ. Its meaning depends on location, size, appearance, growth, and symptoms; some nodules are harmless, while others need testing or follow-up.

\medterm{Noise Induced Hearing Loss And Its Prevention} Noise-induced hearing loss is permanent injury to inner-ear hair cells from repeated or intense sound exposure, often with tinnitus or difficulty understanding speech. Limiting exposure, increasing distance, taking quiet breaks, and using well-fitted hearing protection reduce risk.

\medterm{Noise-Related Hearing Loss} Alternate name; see \textit{Hearing Loss}.

\medterm{Noma} Noma is a rapidly progressive gangrenous infection of the mouth and face, usually in severely malnourished or immunocompromised children. It destroys oral and facial tissue and is treated with antibiotics, nutritional support, and sometimes surgery.

\medterm{Nomenclature} Nomenclature is the organized system used to name and classify diseases, organisms, chemicals, anatomy, medicines, or other entities.

\medterm{Nomogram} A graphical tool that estimates a quantity or risk from related measurements. The Rumack–Matthew nomogram is used in selected cases of acute acetaminophen ingestion.

\medterm{Nonaccidental Injury} An injury caused deliberately or by neglect rather than by an accident. In children, suspicious patterns may include injuries inconsistent with the history, injuries at different stages of healing, or injuries in protected areas.

\medterm{Non-Accidental Trauma} An injury caused deliberately or through abuse rather than by an accident. Clinicians consider it when the explanation does not fit the injuries, when injuries are repeated or occur at different healing stages, or when a child cannot explain what happened; suspected cases require safeguarding assessment.

\medterm{Non-Adherence} Failure to follow an agreed treatment plan, which may result from access barriers, side effects, beliefs, misunderstanding, or other circumstances.

\medterm{Non-Alcoholic Fatty Liver Disease} A spectrum of liver disorders characterized by hepatic steatosis exceeding 5 percent of hepatocytes in individuals without significant alcohol consumption, ranging from simple steatosis to non-alcoholic steatohepatitis (NASH), cirrhosis, and hepatocellular carcinoma.

\synonyms
NAFLD; Metabolic Dysfunction-Associated Steatotic Liver Disease; MASLD

\medterm{Nonalcoholic Fatty Liver Disease} The former name for metabolic dysfunction-associated steatotic liver disease (MASLD), in which excess fat builds up in the liver together with at least one cardiometabolic risk factor, such as obesity, diabetes, high blood pressure, or abnormal cholesterol. The older name remains common in medical records.

\medterm{Non-Alcoholic Steatohepatitis} An aggressive, progressive form of non-alcoholic fatty liver disease (NAFLD/MASLD) characterized by hepatic steatosis accompanied by lobular necroinflammation and hepatocyte ballooning degeneration, with or without Mallory-Denk bodies and pericellular fibrosis, predisposing to cirrhosis and hepatocellular carcinoma.
\synonyms
NASH; Metabolic Dysfunction-Associated Steatohepatitis (MASH)

\medterm{Non-Allergic Rhinitis} A chronic nasal syndrome characterized by sneezing, rhinorrhea, nasal congestion, and post-nasal drip without clinical or laboratory evidence of IgE-mediated allergic sensitization, frequently triggered by weather changes, cold dry air, chemical odors, or foods (vasomotor rhinitis).

\synonyms
Vasomotor Rhinitis; NAR

\medterm{Nonallergic Rhinitis} Nasal congestion, sneezing, or runny nose caused by irritants or other triggers rather than an allergic immune response. Triggers may include changes in temperature or humidity, smoke, odors, medicines, or hormonal changes.

\synonyms
Vasomotor Rhinitis; Idiopathic Rhinitis

\medterm{Nonallergic Rhinopathy} Nonallergic rhinopathy causes congestion, runny nose, or sneezing without an allergic immune reaction, often triggered by odors, smoke, weather changes, or irritants.

\medterm{Non-Articular Rheumatism} An older term for painful disorders of muscles, tendons, bursae, or other tissues outside the joints.

\medterm{Non-Asthmatic Eosinophilic Bronchitis} A cause of long-lasting cough in which certain white blood cells inflame the airways without the variable narrowing seen in asthma. Breathing tests may be normal, but an inhaled corticosteroid can reduce the cough when this diagnosis is confirmed.

\medterm{Nonbacterial Prostatitis} Chronic pelvic pain syndrome (CP/CPPS) characterized by persistent pelvic, perineal, or genital pain and urinary symptoms for at least 3 months without evidence of active urinary tract bacterial infection on standard cultures.

\synonyms
Chronic Pelvic Pain Syndrome in Men; Abacterial Prostatitis; Category III Prostatitis

\medterm{Nonbacterial Thrombotic Endocarditis} Small, sterile blood clots that form on heart valves without infection. It is often linked to cancer, lupus, or another condition that makes blood clot easily. The clots can break off and cause stroke or blockage elsewhere, so treatment focuses on the cause and preventing further clots.

\synonyms
NBTE; Marantic endocarditis; Libman-Sacks-like nonbacterial endocarditis

\medterm{Non-Celiac Gluten Intolerance} Alternate name; see \textit{Non-Celiac Gluten Sensitivity}.

\medterm{Non-Celiac Gluten Sensitivity} Symptoms attributed to gluten or wheat that improve after removal, despite testing not showing celiac disease or wheat allergy. Diagnosis requires careful medical evaluation because other foods and conditions can cause similar symptoms.

\medterm{Non-Communicable Disease} A chronic, non-transmissible medical condition that is not caused by acute infectious agents and typically progresses slowly over a prolonged duration. Arising from a multifactorial interplay of genetic, physiological, environmental, and behavioral factors, the four primary groups responsible for the vast majority of global mortality and disability are cardiovascular diseases (such as coronary artery disease, stroke, and hypertension), cancers, chronic respiratory diseases (such as COPD and asthma), and diabetes mellitus. Management and public-health prevention focus on early screening, long-term pharmacological control, and modifying shared behavioral risk factors, including tobacco use, physical inactivity, unhealthy diet, and harmful alcohol consumption.

\synonyms
NCD; Non-Communicable Diseases; Noncommunicable Disease; Chronic Non-Communicable Disease

\medterm{Non Communicating Hydrocephalus} A buildup of cerebrospinal fluid in the brain because a blockage prevents it from moving between the brain’s fluid-filled spaces and reaching the area where it is absorbed. It can increase pressure on the brain and cause headache, vomiting, sleepiness, vision problems, or difficulty walking.

\medterm{Noncompetitive Antagonist} A medicine that reduces another medicine or natural chemical's effect by attaching to its receptor or signaling pathway in a way that extra amounts of the activating substance cannot fully overcome.

\medterm{Noncompetitive Inhibition} A form of enzyme inhibition in which a substance binds away from the enzyme’s working site and reduces its activity. Because it does not compete directly with the usual substrate, adding more substrate does not fully overcome the inhibition.

\medterm{Non-Compliant Angioplasty Balloon} A small balloon on a catheter that keeps its shape when inflated at high pressure to widen a narrowed blood vessel. It is used during selected vascular procedures and can rarely cause bleeding, vessel injury, or blockage.
\synonyms
High-Pressure Angioplasty Balloon; Non-Compliant Balloon

\medterm{Non-Contact Tonometry} A quick estimate of pressure inside the eye using a brief puff of air against the clear front surface of the eye. It does not touch the eye and can screen for glaucoma, but unusual results may need confirmation.

\medterm{Non-Convulsive Status Epilepticus} A prolonged state of continuous electrographic seizure activity lasting 30 minutes or longer characterized by altered mental status, confusion, coma, or subtle motor automatisms (such as eyelid fluttering) in the complete absence of generalized tonic-clonic motor convulsions, diagnosed via continuous EEG.

\synonyms
NCSE

\medterm{Nonconvulsive Status Epilepticus} Ongoing seizure activity without prominent rhythmic limb jerking, often presenting as prolonged confusion, impaired awareness, or subtle twitching. An electroencephalogram can confirm it, and treatment is urgent.

\medterm{Nondisjunction} Failure of chromosome copies to separate properly when a cell divides. This can produce cells with too many or too few chromosomes and may cause conditions such as Down syndrome, Turner syndrome, or some miscarriages.

\medterm{Non-Fatal Drowning} Survival after breathing is impaired by submersion or immersion in liquid. A person may still develop lung injury, brain injury, or delayed breathing problems, even after apparently recovering.

\medterm{Non-Functioning Pituitary Adenoma} A usually noncancerous pituitary tumor that does not release enough active hormones to cause a hormone syndrome. As it enlarges, it may cause headache, reduced vision, or loss of normal pituitary function.

\medterm{Non-Hodgkin Lymphoma} A group of cancers that begin in immune cells called lymphocytes. Some grow slowly and others quickly; they may cause swollen lymph nodes, fever, night sweats, weight loss, or symptoms in an affected organ.

\synonyms
NHL; B-Cell or T-Cell Lymphoma

\medterm{Non-Hodgkin Lymphoma In Children} A group of childhood cancers arising from immune cells called lymphocytes. Some grow quickly and can cause swollen nodes, fever, abdominal symptoms, or breathing problems; the exact type and spread determine treatment.

\medterm{Non-Hodgkin's Lymphoma} Alternate spelling; see \textit{Non-Hodgkin lymphoma}.

\medterm{Non-Hodgkins Lymphomas} Alternate spelling and plural form; see \textit{Non-Hodgkin lymphoma}.

\medterm{Non-Infectious Cystitis} Inflammation of the bladder without infection by bacteria or another microorganism. It may cause pain, urgency, and frequent urination after irritants, radiation, medicines, chemical exposure, or other bladder conditions.

\medterm{Noninfectious Cystitis} A heterogeneous group of bladder inflammatory conditions occurring in the absence of active bacterial infection, encompassing radiation cystitis, chemotherapy-induced cystitis (e.g., cyclophosphamide), chemical cystitis, and interstitial cystitis.

\synonyms
Abacterial Cystitis; Sterile Cystitis; Chemical/Radiation Cystitis

\medterm{Noninvasive} Not involving a surgical cut, a procedure entering the body, or a device placed inside it. Examples include measuring blood pressure with a cuff, external imaging, and breathing support through a mask. Risks depend on the method used.

\medterm{Non-Invasive Positive Pressure Ventilation} The delivery of mechanical ventilatory support without an invasive endotracheal tube or tracheostomy, utilizing a tight-fitting facial or nasal mask to provide continuous positive airway pressure (CPAP) or bilevel positive airway pressure (BiPAP) in acute hypercapnic respiratory failure or cardiogenic pulmonary edema.

\synonyms
NIPPV; NPPV

\medterm{Non-Invasive Prenatal Testing} A blood test during pregnancy that estimates the chance of selected chromosome conditions, especially trisomy 21, 18, and 13, by analyzing small DNA fragments from the placenta. It is screening, not diagnosis; an abnormal result needs confirmatory testing.

\medterm{Non-Invasive Procedure} An examination or treatment that does not cut or pierce the skin or place an instrument inside the body. Examples include blood-pressure checks, ECGs, ultrasound, MRI, X-rays, and physical therapy. These procedures are generally safer than invasive procedures, but each can still have limits or risks.

\synonyms
Noninvasive Procedure; Non-Invasive Medical Procedure; Noninvasive Test; Non-Invasive Technique

\medterm{Non-Invasive Ventilation} Breathing support delivered through a face or nose mask instead of a tube placed in the windpipe. Continuous or two-level air pressure can improve breathing in selected people with respiratory failure, sleep apnea, or fluid in the lungs.

\medterm{Nonischemic Priapism} Alternate name; see \textit{Priapism}.

\medterm{Nonketotic Hyperglycemic Hyperosmolar Syndrome} A life-threatening diabetic emergency in which very high blood sugar causes profound dehydration and concentrated blood, usually without major ketone buildup. It can cause extreme thirst, weakness, confusion, seizures, coma, and organ failure.

\medterm{Non-Ketotic Hyperosmolar State} A life-threatening acute metabolic emergency occurring predominantly in type 2 diabetes mellitus, characterized by extreme hyperglycemia (glucose greater than 600 mg/dL), profound hyperosmolality (plasma effective osmolality greater than 320 mOsm/kg), and severe dehydration in the absence of significant ketoacidosis.
\synonyms
Hyperosmolar Hyperglycemic State; HHS

\medterm{Nonmelanoma Skin Cancer} A group of skin cancers that arise from the cells of the epidermis, excluding melanoma.
The two most common types are basal cell carcinoma (BCC) and squamous cell carcinoma
(SCC), which together account for the vast majority of all skin cancers diagnosed
globally.

\medterm{Non-Muscle-Invasive Bladder Cancer} Bladder cancer limited to the inner lining or connective tissue beneath it, without invasion of the bladder muscle. It can recur and may progress to muscle-invasive disease.

\medterm{Non-Obstructive Azoospermia} Complete absence of sperm in semen because the testicles make too few or no sperm, rather than because a tube is blocked. Causes include genetic conditions, undescended testes, infections, cancer treatment, or an unknown problem. Testing may include hormone tests and genetic studies; some people can use surgically retrieved sperm for assisted reproduction.

\synonyms
NOA; Secretory Azoospermia; Testicular Failure Azoospermia

\medterm{Nonpathogenic} Nonpathogenic means not ordinarily capable of causing disease in a given host or setting; an organism can still cause opportunistic infection when immunity or barriers are impaired.

\medterm{Non-Pitting Edema} Swelling of subcutaneous tissue that does not leave a persistent indentation when firm digital pressure is applied, usually due to local accumulation of protein-rich fluid, mucopolysaccharides, or chronic cellular infiltration. Common causes include lymphedema, pretibial myxedema in Graves disease, and severe lipedema.

\medterm{Non-Polio Enterovirus Infection} Alternate name; see \textit{Enterovirus}.

\medterm{Non-Proliferative Diabetic Retinopathy} An early form of diabetes-related damage to the retina, the light-sensitive tissue at the back of the eye. Small blood vessels may leak or swell, causing blurred vision, but abnormal new vessels have not formed.

\medterm{Nonproliferative Retinopathy} Retinal disease in which small blood vessels leak or become damaged without abnormal new vessels growing. Diabetes is a common cause; progression can impair vision, especially when the macula is affected.
\synonyms
Background Retinopathy; NPDR

\medterm{Non-Radiographic Axial Spondyloarthritis} Inflammatory arthritis affecting the spine or joints between the spine and pelvis without definite changes on a plain X-ray. It may cause persistent back pain and morning stiffness, and may occur with eye inflammation, psoriasis, or inflammatory bowel disease.

\medterm{Non-Rapid Eye Movement Sleep} Sleep that includes lighter and deeper stages without rapid eye movements. Deep stages support physical restoration, while lighter stages help the body transition into and out of sleep.

\medterm{Non-Rebreather Mask} A face mask with a reservoir bag and one-way valves that can deliver a high concentration of supplemental oxygen when fitted and used correctly.

\medterm{Non-Rebreather Mask} A high-flow oxygen delivery device featuring an attached reservoir bag and one-way directional valves that prevent entrainment of room air and rebreathing of exhaled gas, capable of delivering inspired oxygen fractions (FiO2) between 80 and 95 percent at flow rates of 10 to 15 L/min.

\synonyms
NRB Mask

\medterm{Non-REM Sleep} Sleep stages without rapid eye movements, ranging from light sleep to deep slow-wave sleep. These stages support physical restoration, memory processing, and the normal transition between waking and sleep.
\synonyms
NREM sleep

\medterm{Non-Responder} A person or disease that shows little or no expected response to a treatment, vaccine, or test stimulus.

\medterm{Nonselective Beta-Blockers} Medicines that block two main types of adrenaline receptor. They slow the heart and lower blood pressure but can also narrow the airways, so people with asthma or certain lung diseases may need a different medicine.

\medterm{Nonseminomatous Germ Cell Tumor} A cancer arising from reproductive cells, most often in the testicle. It may grow and spread quickly; blood tests, scans, and tissue examination guide staging and treatment.

\synonyms
NSGCT; Nonseminoma

\medterm{Non-Seminomatous Germ Cell Tumor} A major clinicopathological category of testicular and extragonadal germ cell neoplasms comprising embryonal carcinoma, yolk sac tumor, choriocarcinoma, and teratoma (or mixed variants). Characterized by aggressive clinical behavior with early hematogenous and retroperitoneal lymphatic dissemination, producing characteristic serum tumor markers (alpha-fetoprotein and human chorionic gonadotropin); treated with radical orchiectomy followed by platinum-based combination chemotherapy.

\synonyms
NSGCT; Nonseminoma

\medterm{Nonsense Variant} A change in DNA that tells a cell to stop making a protein too early. The resulting shortened or missing protein may cause disease, but the effect depends on the gene, the location of the change, and how much normal protein remains.

\medterm{Non-Small Cell Lung Cancer} The most common group of lung cancers, including adenocarcinoma and squamous-cell cancer. It may start in different parts of the lung and can spread; treatment depends on its stage, cell type, gene changes, and the person's health.

\medterm{Non-Small Cell Lung Carcinoma} A broad histological grouping of primary epithelial malignancies of the lung comprising approximately 85 percent of all lung cancers, including adenocarcinoma, squamous cell carcinoma, and large cell carcinoma, distinguished by clinical behavior and targeted therapy sensitivity from small cell lung cancer.

\synonyms
NSCLC

\medterm{Nonspecific Interstitial Pneumonia} A pattern of lung inflammation or scarring that develops relatively evenly over time. It is often linked to connective-tissue disease, medicines, or an unknown cause and can cause breathlessness and a dry cough.

\medterm{Nonspecific Interstitial Pneumonitis} A pattern of lung inflammation and scarring that develops fairly evenly through the lung tissue. It may cause a dry cough and breathlessness and is often linked to autoimmune disease, medicines, or no identified cause.

\medterm{Non-Specific Urethritis} Inflammation of the tube that carries urine out of the body when a specific cause has not yet been identified. It may cause burning urination or discharge and is often linked to sexually transmitted infections, so testing and treatment of partners may be needed.

\medterm{Non-ST Elevation Myocardial Infarction} An acute coronary syndrome characterized by ischemic myocardial necrosis with detectable elevated cardiac biomarkers (cardiac troponin I or T) in the absence of acute persistent ST-segment elevation on ECG (presenting instead with ST-segment depression, T-wave inversion, or normal tracing).
\synonyms
NSTEMI; Non-ST-Segment Elevation Infarct

\medterm{Non-ST-Elevation Myocardial Infarction} A heart attack in which blood tests show heart-muscle injury but the electrocardiogram does not show persistent ST-segment elevation. It is usually caused by reduced or blocked coronary blood flow and may require urgent assessment and treatment.

\medterm{Nonsteroidal Anti-Inflammatory Drugs} A group of medicines that reduce pain, fever, and inflammation by lowering the body’s production of prostaglandins. They can cause stomach bleeding, kidney injury, fluid retention, heart problems, or allergic reactions, especially at high doses or with long use.

\medterm{Nonsteroidal Anti-Inflammatory Drugs And Ulcers} Nonsteroidal anti-inflammatory drugs can reduce stomach-protective prostaglandins and cause gastritis, ulcers, and gastrointestinal bleeding. Risk rises with age, prior ulcer, anticoagulants, steroids, kidney disease, and high doses; safer alternatives or acid protection may be needed.

\synonyms
NSAID-Induced Peptic Ulcer; NSAID Gastropathy

\medterm{Non-Stress Test} A test during pregnancy that monitors the baby's heart rate while the baby moves. A healthy response usually includes temporary heart-rate increases; an abnormal result may reflect sleep, medicines, or a problem requiring further assessment.

\medterm{Non-ST-Segment Elevation Myocardial Infarction} An acute coronary syndrome characterized by evidence of myocardial necrosis (elevated cardiac troponin biomarkers) in the absence of acute ST-segment elevation on electrocardiogram, typically reflecting subendocardial ischemia caused by a non-occlusive mural thrombus.

\synonyms
NSTEMI

\medterm{Nonsyndromic} Nonsyndromic means occurring without the additional characteristic features of a recognized syndrome; it may still be genetic, congenital, or clinically significant.

\medterm{Nontropical Sprue} Alternate name; see \textit{Celiac Disease}.

\medterm{Nontuberculous Mycobacteria} Mycobacteria other than the tuberculosis complex and M. leprae. They are found in soil and water and can cause lung, skin, lymph-node, or disseminated infections, especially in people with underlying illness.

\medterm{Nontuberculous Mycobacterial Infection} Infection caused by Mycobacterium species other than M. tuberculosis, often affecting the lungs but sometimes the lymph nodes, skin, or other organs. Disease is more likely in people with structural lung disease or impaired immunity.

\medterm{Nontuberculous Mycobacterial Lung Disease} Chronic lung infection caused by environmental Mycobacterium species other than M. tuberculosis.
It may produce cough, fatigue, weight loss, breathlessness, or coughing blood, especially in people with structural
lung disease. Diagnosis requires compatible imaging and repeated microbiologic testing; treatment is prolonged and individualized.

\medterm{Nonulcer Dyspepsia} Alternate name; see \textit{Functional Dyspepsia}.

\medterm{Nonulcer Stomach Pain} Alternate name; see \textit{Functional Dyspepsia}.

\medterm{Nonunion And Malunion} Nonunion means a broken bone fails to heal after an expected period, leaving pain or instability. Malunion means it heals in a poor position, causing deformity, altered movement, or strain on nearby joints. Both may require further treatment.

\medterm{Nonunion Fracture} A broken bone that has failed to heal after the expected period and shows no further progress toward joining. It may cause ongoing pain, movement at the fracture, weakness, or deformity and may require surgery or other treatment.

\medterm{Noonan Syndrome} An inherited condition that may cause distinctive facial features, short stature, heart defects, delayed development, easy bruising, or swelling from lymph problems. Its features vary, and regular heart, growth, hearing, and developmental checks are important.

\medterm{Noonan Syndrome With Multiple Lentigines} A genetic disorder also called LEOPARD syndrome, characterized by multiple lentigines, distinctive facial features, heart disease, short stature, and sometimes deafness or genital abnormalities.

\medterm{Noradrenaline} A chemical messenger acting as both a hormone and neurotransmitter. It helps regulate alertness, attention, stress responses, heart activity, and blood-vessel tightening, which can raise blood pressure.

\synonyms
Norepinephrine

\medterm{No-Reflow Phenomenon} Failure of blood to reach small vessels and damaged tissue even after a larger blocked artery has been reopened. It can occur after a heart attack or stroke and may limit recovery despite successful treatment of the main blockage.

\medterm{Norepinephrine} A hormone and nerve messenger involved in alertness, stress responses, heart function, and blood-vessel tightening. It can increase blood pressure and is also used in hospitals to support blood pressure during some forms of shock.
\synonyms
Noradrenaline

\medterm{Norethisterone} A manufactured progesterone-like medicine used for contraception and selected menstrual problems or endometriosis. It can change bleeding patterns and may cause nausea, headache, mood changes, breast tenderness, or other hormone-related effects.

\medterm{Norgestrel} A synthetic progestogen used in some hormonal contraceptives. It helps prevent pregnancy by suppressing ovulation and changing cervical mucus and the uterine lining.

\medterm{Norketamine} Norketamine is an active metabolite of ketamine with NMDA-receptor antagonist activity; it contributes to ketamine's analgesic and psychoactive effects and may persist after the parent drug is cleared.

\medterm{Normal} Within an expected range or consistent with usual anatomy, physiology, or function for a specified population.

\medterm{Normal Flora} The bacteria, fungi, and other microorganisms that normally live on the skin and in the mouth, intestine, and urinary or genital tract without causing illness. They can help block harmful germs, support immunity, and contribute to digestion and vitamin production.

\medterm{Normal Growth And Development} Normal growth and development describes expected physical growth, movement, language, learning, and social progress for a person’s age and stage of life.

\medterm{Normalization} A mathematical adjustment that makes measurements comparable by correcting for known differences in scale, background, equipment response, or other measurement conditions.

\medterm{Normal Labour} Spontaneous labour at term with regular uterine contractions, progressive cervical dilatation, and delivery through the birth canal without major complication.

\medterm{Normal Pressure Hydrocephalus} A buildup of cerebrospinal fluid that enlarges spaces inside the brain and may cause a slow, unsteady walk, reduced thinking, and loss of bladder control. Some people improve with specialist evaluation and drainage treatment.

\medterm{Normal-Pressure Hydrocephalus} A disorder in which cerebrospinal fluid accumulates and causes gait difficulty, cognitive decline, and urinary problems despite near-normal measured pressure.

\medterm{Normal Range} A normal range is the interval of results commonly found in a reference population; a result outside it does not by itself prove disease.

\medterm{Normal Saline} A sterile approximately 0.9\% sodium-chloride solution used for fluid replacement, irrigation, and dilution of medicines.

\medterm{Normal-Tension Glaucoma} A type of open-angle glaucoma in which the optic nerve is damaged and vision may be lost even though pressure inside the eye remains within the normal range.

\medterm{Normoblast} A nucleated immature red-cell precursor in bone marrow that normally loses its nucleus before entering circulation.

\medterm{Normocalcemic Hyperparathyroidism} Condition of elevated parathyroid hormone with persistently normal serum calcium and
normal vitamin D levels. Exclusion of secondary causes of elevated PTH including vitamin
D deficiency, renal insufficiency, and medications is required. Represents an early or
mild form of primary hyperparathyroidism. May progress to overt hypercalcemia.

\medterm{Normocytic Anaemia} Alternate British spelling; see \textit{Normocytic Anemia}.

\medterm{Normocytic Anemia} A form of anemia in which circulating red blood cells are of normal average size (mean corpuscular volume [MCV] between $80$ and $100\,\text{fL}$), but total red blood cell count and hemoglobin concentration are decreased. Diagnostic evaluation relies on the reticulocyte count to differentiate between: (1) hyperproliferative causes with elevated reticulocytes, such as acute blood loss (hemorrhage) or premature red cell destruction (hemolytic anemia), and (2) hypoproliferative causes with low reticulocytes, including anemia of chronic disease, chronic kidney disease (erythropoietin deficiency), aplastic anemia, and bone marrow infiltration. Symptoms include fatigue, weakness, dizziness, and pallor.
\synonyms{Normocytic Normochromic Anemia; Normal-MCV Anemia}

\medterm{Normotensive} Having blood pressure within the expected range for the individual and measurement setting. Having blood pressure within the expected range for a person, based on the measurement setting.

\medterm{Norovirus} A highly contagious virus that causes acute gastroenteritis, or inflammation of the stomach and intestines. Symptoms usually begin 12--48 hours after exposure and include sudden nausea, vomiting, watery diarrhea, stomach cramps, and sometimes fever, headache, or body aches. It spreads mainly when tiny particles of stool or vomit enter the mouth, through infected people, contaminated food or water, surfaces, or particles released during vomiting. Most illness improves within 1--3 days, but severe fluid loss can cause dehydration, especially in infants, older adults, and people with other illnesses.

\medterm{Norovirus Infection} Acute gastroenteritis caused by norovirus. Symptoms usually begin 12--48 hours after exposure and include nausea, vomiting, watery diarrhea, abdominal cramps, and sometimes fever, headache, or body aches. Infection spreads easily through infected people, contaminated food or water, surfaces, or particles released during vomiting. Most cases improve within 1--3 days, but fluid loss may cause dehydration, especially in infants, older adults, and people with other illnesses.

\medterm{Norrie Disease} A rare inherited eye disorder in which the retina develops abnormally, usually causing severe vision loss or blindness from birth or early infancy. Some people also develop hearing loss, developmental differences, or problems in the eye’s blood vessels.

\medterm{Nortriptyline} A medicine used to treat depression and sometimes long-lasting nerve pain. It changes the levels of certain chemical messengers in the nervous system and may cause dry mouth, constipation, sleepiness, dizziness, heart-rhythm problems, or overdose toxicity.

\medterm{Norwegian Scabies} A severe form of scabies in which thick crusts contain very large numbers of mites. It is extremely contagious and occurs especially in people with weakened immunity or difficulty sensing or scratching the skin, requiring prompt treatment and contact precautions.

\medterm{Norwood Procedure} The first operation for many newborns with hypoplastic left heart syndrome and related heart defects. It rebuilds the main artery so the right ventricle can send blood to the body, creates a wide opening between the upper heart chambers, and provides controlled blood flow to the lungs through a shunt or tube. It is usually performed within the first days of life and is followed by further staged operations.

\synonyms
Stage 1 Norwood Reconstruction; Norwood Operation

\medterm{Nose} The external structure and internal passages used for breathing and smell. The nose warms, moistens, and filters air and contains bone, cartilage, skin, mucus, blood vessels, and smell-sensing tissue.

\medterm{Nosebleed} Singular form; see \textit{Nosebleeds}.

\medterm{Nosebleeds} Bleeding from blood vessels inside the nose, commonly caused by dryness, irritation, injury, infection, allergies, medicines, or high blood pressure. Episodes may be brief or recurrent and can vary in amount.

\medterm{Nose Fracture} A break in one or more nasal bones caused by trauma, potentially causing pain, swelling, deformity, bleeding, or nasal obstruction.

\medterm{Nosocomial} Acquired during healthcare, especially in a hospital, rather than being present or incubating when the person entered care. A nosocomial infection may be caused by bacteria, viruses, fungi, or other organisms and can involve drug-resistant organisms.

\medterm{Nosocomial Infection} A hospital-acquired infection developing during or after hospitalization that was not
present or incubating at admission, classically defined by onset beyond 48 hours. Common
sites are bloodstream, urinary tract, lungs, and surgical wounds; causative organisms
are frequently multidrug-resistant, prevented by hand hygiene and stewardship.

\medterm{Nosocomial Norovirus Infection} Healthcare-associated norovirus gastroenteritis characterized by acute, explosive onset of watery diarrhea, severe non-bloody emesis, and abdominal cramps. Outbreak management necessitates strict patient isolation, contact precautions, rigorous hand hygiene with soap and water (as alcohol-based rubs are less effective against non-enveloped virions), and hypochlorite-based surface disinfection.

\synonyms
Hospital-Acquired Norovirus; Healthcare Norovirus Outbreak; Hospital Winter Vomiting Disease

\medterm{Nosocomial Vancomycin-Resistant Enterococci} Enterococcal bacteria that resist vancomycin and can spread in healthcare settings. They may cause urinary, wound, bloodstream, or other infections, particularly in people who are seriously ill or have devices such as catheters; treatment requires susceptibility testing.

\medterm{Nostril} One of the two openings at the front of the nose. Each nostril leads into the nasal passage, where hairs and mucus help trap particles before air reaches the deeper airways.

\medterm{Notencephalocele} A rare congenital neural-tube defect involving protrusion of brain tissue or meninges through a skull defect; terminology varies by source.

\medterm{Notifiable Disease} An infection or other condition that healthcare professionals or laboratories must report to public-health authorities under local law. Reporting helps detect outbreaks, arrange control measures, protect contacts, and plan services; the exact list and reporting rules vary by location.

\medterm{Notochord} A flexible embryonic rod that guides axial development and largely disappears as the vertebral column forms.

\medterm{NOWS} Alternate name; see \textit{Neonatal Opioid Withdrawal Syndrome}.

\medterm{Noxious} Harmful, poisonous, or capable of causing tissue injury or pain. Noxious stimuli activate protective sensory pathways such as pain and withdrawal reflexes.

\medterm{NREM Sleep} The part of sleep without rapid eye movements, ranging from light sleep to deep slow-wave sleep. It alternates with REM sleep in repeated cycles and supports physical restoration, learning, and memory.

\medterm{NSAID} A nonsteroidal anti-inflammatory drug used to reduce pain, fever, and inflammation. Examples include ibuprofen and naproxen; these medicines can irritate the stomach, cause bleeding or kidney problems, and increase heart risks in some people.

\synonyms
Nonsteroidal anti-inflammatory drug

\medterm{NSAID-Induced Ulcers} Alternate name; see \textit{Peptic Ulcer Disease}.

\medterm{NSAIDs} A class of analgesic, antipyretic, and anti-inflammatory drugs that inhibit
cyclooxygenase (COX-1 and/or COX-2), reducing prostaglandin and thromboxane synthesis.

\synonyms
Non-Steroidal Anti-Inflammatories; COX Inhibitors

\medterm{NSTEMI} A heart attack in which a blocked or narrowed coronary artery injures heart muscle but does not produce the typical persistent ST-segment elevation on an ECG. It may cause chest pressure, breathlessness, sweating, nausea, or faintness.

\synonyms
Non-ST-Elevation Myocardial Infarction

\medterm{NTD} NTD commonly means neural-tube defect, a congenital brain or spinal-cord malformation associated with inadequate folate and other factors.

\medterm{Nuchal} Relating to the back of the neck. Nuchal rigidity means painful or limited neck bending and can occur with meningitis or bleeding around the brain. Nuchal translucency is a first-trimester ultrasound measurement used to estimate certain pregnancy risks.

\medterm{Nuchal Cord} Wrapping of the umbilical cord around the fetal neck, occurring once or more times. It is common and often harmless, but monitoring during labor is needed if fetal heart-rate abnormalities occur.

\medterm{Nuchal Rigidity} Stiffness or pain that makes it difficult to bend the neck forward. It may occur with meningitis, bleeding around the brain, severe muscle problems, or other illnesses; sudden stiffness with fever or severe headache may signal a serious condition.

\medterm{Nuchal Translucency Screening} An ultrasound measurement of the clear space at the back of a fetus’s neck, usually performed early in pregnancy. A larger-than-expected measurement can indicate increased risk of chromosome or heart problems but is not a diagnosis.

\medterm{Nuchal Translucency Test} An early-pregnancy ultrasound measurement of the clear space at the back of the fetus’s neck. It estimates the chance of certain chromosome or structural problems but cannot diagnose them; abnormal results need follow-up testing.

\medterm{Nuclear} Relating to the nucleus of a cell, the center of an atom, or medical tests and treatments that use radioactivity. The intended meaning depends on the surrounding medical context.

\medterm{Nuclear Factor Kappa B} A ubiquitous inducible master transcription factor complex that coordinates immune and inflammatory responses, cell survival, and cytokine transcription; normally held inactive in the cytoplasm by I-kappa-B inhibitor proteins until activated by inflammatory cytokines, bacterial LPS, or viral infection.

\synonyms
NF-kappaB; NF-kB

\medterm{Nuclear Medicine} A medical specialty that uses small amounts of radioactive substances to diagnose or treat disease. Scans show body function or chemical activity, while treatments deliver radiation to selected tissues.

\medterm{Nuclear Medicine Imaging} Imaging that uses a radioactive tracer and a specialized detector to show organ function, blood flow, chemical activity, or disease distribution. Common examples include PET, SPECT, bone scans, thyroid scans, and radionuclide cardiac imaging.

\medterm{Nuclear Pore Complex} A large, multiprotein cylindrical channel spanning the double-membrane nuclear envelope, composed of multiple nucleoporin proteins that regulate and mediate the selective bidirectional transport of macromolecules (RNA, ribonucleoproteins, proteins) between the nucleus and the cytoplasm.

\synonyms
NPC

\medterm{Nuclear Sclerosis} A common age-related cataract in which the center of the eye’s lens gradually becomes yellow and hard. It can slowly blur vision, reduce contrast, or alter color perception and is treated when it interferes with daily activities.

\medterm{Nuclear Stress Test} A test that compares blood flow to heart muscle at rest and during exercise or medicine-induced stress using a radioactive tracer. It can identify areas of reduced perfusion that suggest coronary artery disease.

\synonyms
Nuclear medicine tests

\medterm{Nuclear Ventriculography} A heart scan that follows a small radioactive tracer through the bloodstream to measure how much blood the heart pumps with each beat and whether the heart walls move normally.

\medterm{Nucleic Acid} DNA or RNA, the molecules that store, copy, and use genetic instructions. DNA usually stores the long-term information, while RNA helps read those instructions and make proteins or perform other jobs in cells.

\medterm{Nucleic Acid Hybridization} A laboratory method in which matching DNA or RNA strands bind together. This allows a test to find, identify, or measure a particular genetic sequence from a person, tumor, or infectious organism.

\medterm{Nucleolus} A dense region within the cell nucleus where ribosomal RNA is produced and ribosome subunits begin assembly.

\medterm{Nucleoside} A small molecule made of a nitrogen-containing base joined to a sugar. Adding one or more phosphate groups makes a nucleotide, the building block used to make DNA or RNA and to carry energy in cells.

\medterm{Nucleoside Analogs} Medicines that resemble the building blocks used to make DNA or RNA. Some viruses use them by mistake, which interrupts viral copying; the medicine chosen depends on the virus, illness, kidney function, and other factors.

\medterm{Nucleotide} A small molecule made of a sugar, a nitrogen-containing base, and one or more phosphate groups. Nucleotides build DNA and RNA and also carry energy used by cells.

\medterm{Nucleus} The membrane-bound cell organelle containing most genetic material, or a defined cluster of neurons in the central nervous system.

\medterm{Nucleus Accumbens} A group of nerve cells deep in the brain that helps process reward, motivation, learning, and reinforcement. Its activity is involved in normal behavior and in disorders such as addiction and some psychiatric illnesses.

\medterm{Nucleus Ambiguus} A longitudinal column of branchial motor lower motor neurons located within the medullary reticular formation that provides motor innervation to the skeletal striated muscles of the soft palate, pharynx, and larynx via cranial nerves IX (glossopharyngeal), X (vagus), and cranial XI.

\medterm{Nucleus Pulposus} The soft, gel-like center of a spinal disk, surrounded by a tough outer ring. It helps spread pressure between vertebrae, but a tear may allow it to bulge or press on nerves.

\medterm{Nulligravida} An obstetric term for a person who has never been pregnant. It is different from nullipara, which means never having given birth beyond the stage of pregnancy used in a particular medical record system.

\medterm{Nullipara} A person who has never carried a pregnancy beyond the stage of fetal viability; the adjective is nulliparous.

\medterm{Null Mutation} A genetic change that prevents a gene from producing a functional protein or RNA. The effect depends on the gene, the other copy, and whether the change occurs in body cells or reproductive cells.

\medterm{Numb} Numb means lacking normal sensation or emotional responsiveness; physical numbness can result from nerve, brain, circulation, metabolic, or medication-related problems.

\medterm{Number Needed To Harm} The estimated number of people who must receive a treatment, rather than the comparison treatment, for one extra person to experience a specified harmful outcome during a stated period. A smaller number means the harm is more frequent.

\medterm{Number Needed To Treat} The number of people who must receive an intervention for one additional person to
experience the specified beneficial outcome compared with the control over a defined
time. It is the reciprocal of absolute risk reduction.

\medterm{Numbness} Numbness is reduced or absent sensation, often from nerve compression, neuropathy, stroke, spinal disease, circulation problems, or anesthetic effects.

\medterm{Numbness And Tingling} Unusual reduced sensation or prickling, often caused by nerve compression, neuropathy, spinal disease, vitamin deficiency, circulation problems, or other neurologic disorders. Sudden one-sided symptoms, weakness, speech trouble, or facial drooping require emergency assessment.

\synonyms
Paresthesia; Pins and Needles

\medterm{Numbness In Hands} Reduced or absent sensation in one or both hands, commonly caused by nerve compression, neuropathy, injury, circulation problems, or systemic disease.

\medterm{Nummular} Coin-shaped or round. The term is commonly used for nummular eczema, which causes itchy, coin-shaped patches of inflamed skin.

\medterm{Nummular Dermatitis} An inflammatory skin condition causing itchy, coin-shaped patches, often on the arms or legs. Dry skin, irritants, and skin injury may contribute; treatment commonly includes moisturizers and prescribed anti-inflammatory medicines.

\medterm{Nummular Eczema} An inflammatory skin condition causing intensely itchy, coin-shaped red patches, sometimes with blisters, crusting, or scaling, usually on the arms or legs. Dry skin, irritation, or injury can trigger it and similar rashes may need examination.

\medterm{Nursemaid's Elbow} Nursemaid's elbow is subluxation of the radial head after sudden traction on a young child's extended or pronated arm; the child avoids using the arm, and gentle clinician reduction usually restores movement.

\medterm{Nurse Practitioner} A nurse practitioner is an advanced-practice nurse who can assess patients, diagnose and treat many conditions, prescribe medicines where authorized, and provide preventive care.

\medterm{Nursery} A place or service for the care of infants, including specialized newborn care in a hospital.

\medterm{Nursing} The professional practice of caring for people, promoting health, preventing illness, and supporting recovery or comfort.

\medterm{Nursing Home} A residential facility that provides ongoing nursing care, assistance with daily activities, and medical or rehabilitation services for
people who cannot live independently.

\medterm{Nut Allergies} Nut allergy is an immune-mediated reaction to tree nuts or peanuts that can cause hives, gastrointestinal or respiratory symptoms, or life-threatening anaphylaxis.

\medterm{Nut-Allergy} An immune-mediated reaction to one or more nuts, potentially causing urticaria, anaphylaxis, or gastrointestinal and respiratory symptoms.

\medterm{Nutation} A nodding or rocking movement; in anatomy it commonly describes movement of the sacrum relative to the pelvis.

\medterm{Nutcracker Esophagus} An older name for a swallowing disorder in which the esophagus contracts with unusually high pressure. It can cause intermittent chest pain or difficulty swallowing; current testing and clinical assessment distinguish it from reflux and other causes.

\medterm{Nutcracker Syndrome} Compression of the vein draining the left kidney, usually between two nearby arteries. It may cause blood in the urine, pelvic or flank pain, or enlarged veins; mild cases may be monitored, while persistent severe symptoms may need treatment.

\medterm{Nutrient} A substance the body needs for energy, growth, repair, or normal function. Nutrients include protein, fats, carbohydrates, vitamins, minerals, and water; too little or too much of some nutrients can cause illness.

\medterm{Nutrition} How the body takes in, digests, absorbs, and uses food and its nutrients for health.

\medterm{Nutritional Anemia} Anemia caused by inadequate intake, absorption, or use of nutrients needed for red-cell production, especially iron, vitamin B12, or folate.

\medterm{Nutritional Assessment} A structured evaluation of food intake, body measurements, laboratory results, symptoms, and health conditions. It identifies nutritional problems and guides a safe, individualized eating plan or nutrition support.

\medterm{Nutritional Deficiency} A lack of a vitamin, mineral, protein, energy, or other nutrient needed for normal body function. It may result from inadequate intake, poor absorption, increased needs, or disease and can impair growth, blood formation, immunity, nerves, or other organs.

\medterm{Nutritional Dermatoses} Skin changes caused by too little or too much of a nutrient. For example, vitamin C deficiency can cause easy bruising and bleeding gums, while niacin deficiency can cause a sun-sensitive rash, diarrhea, and confusion.

\medterm{Nutritional Disorders} Health problems caused by too little, too much, or poor use of nutrients. Examples include undernutrition, obesity, vitamin or mineral shortages, and problems caused by the body being unable to digest or absorb food properly.

\medterm{Nutritional Screening} A brief check for risk of malnutrition using weight change, food intake, illness, and other factors. A concerning result leads to a fuller assessment and an individualized nutrition plan.

\medterm{Nutrition-Focused Physical Examination} A structured examination for physical findings associated with malnutrition, including
loss of subcutaneous fat, muscle wasting, edema, ascites, reduced functional status, and
micronutrient-related signs. It complements dietary history, anthropometry, laboratory
data, and disease assessment.

\medterm{Nutrition In Pregnancy} Eating enough varied food and obtaining appropriate vitamins and minerals to support the pregnant person and growing fetus. Important nutrients include folate, iron, calcium, vitamin D, iodine, and protein; individual needs and supplements should be reviewed with a prenatal clinician.

\medterm{Nutrition Risk Screening 2002} A hospital screening questionnaire that considers recent food intake, weight change, illness severity, and age to identify adults at risk of malnutrition. A high score prompts a fuller assessment; it is not itself a diagnosis.

\medterm{Nutrition Support} Use of oral nutrition supplements, tube feeding, or intravenous nutrition when a person cannot meet nutritional needs by ordinary eating. The method depends on the digestive tract, illness, nutrition needs, and risks of treatment.

\medterm{Nutritive} Providing nourishment that the body can use for energy, growth, repair, or normal function. Nutritive substances include protein, fat, carbohydrates, vitamins, minerals, and water.

\medterm{Nux Vomica} The seed of Strychnos nux-vomica, containing strychnine and brucine; ingestion can cause severe convulsant poisoning.

\medterm{Nyctalopia} Impaired vision in dim light or darkness, also called night blindness.

\medterm{NYHA Functional Classification} A four-class system describing how heart failure symptoms limit activity: class I causes no limitation, class II limits ordinary activity, class III limits less-than-ordinary activity, and class IV causes symptoms even at rest.

\medterm{Nystagmus} Involuntary rhythmic movement of the eyes. It may occur with inner-ear disease, neurologic disorders, reduced vision, medicines, or normal early development; new nystagmus with dizziness, weakness, or vision change may signal an underlying disorder.

\medterm{Nystagmus Retractorius} An unusual eye movement in which the eyes pull backward and may move toward each other when a person tries to look upward. It is a characteristic sign of a problem in the upper part of the brain, but other findings are needed to identify the cause.
\synonyms
Retraction-Convergence Nystagmus; Koerber-Salus-Elschnig Sign
